% model.tex
% Stand-alone description of the cost-augmented noisy-threshold RSA model.
% Can be embedded in main.tex via % model.tex
% Stand-alone description of the cost-augmented noisy-threshold RSA model.
% Can be embedded in main.tex via % model.tex
% Stand-alone description of the cost-augmented noisy-threshold RSA model.
% Can be embedded in main.tex via % model.tex
% Stand-alone description of the cost-augmented noisy-threshold RSA model.
% Can be embedded in main.tex via \input{model}.
%
% Compile stand-alone:
%   pdflatex model.tex
%
% Embed in main.tex:
%   \input{model}

\documentclass[12pt]{article}
\usepackage{CSPArticleStyle}
\usepackage{csquotes}
\usepackage[margin=1in]{geometry}
\usepackage{booktabs}
\usepackage{parskip}
\usepackage{amsmath}
\usepackage{amssymb}
\usepackage{times}
\usepackage[T1]{fontenc}
\usepackage[utf8]{inputenc}
\usepackage{xcolor}

\title{Cost-Augmented Noisy-Threshold RSA Model for Consensus Markers}
\author{}
\date{}

\begin{document}

%% ============================================================
%% When embedding in main.tex, delete from \documentclass to
%% \begin{document} and delete \end{document} at the bottom.
%% ============================================================

\section{Model}

\subsection{Overview}

We model German consensus markers as lexical items that are licensed over
intervals of a latent utility scale. The model combines three ideas:
\begin{enumerate}
  \item perceived controversy is jointly determined by a propositional and a pragmatic component;
  \item speaker choice depends on a utility trade-off between controversy and persuasive goal strength;
  \item stronger markers incur additional pragmatic costs, which creates a
        \emph{weakest-sufficient} bias and can block overly strong markers.
\end{enumerate}

The three investigated markers are ordered from weakest to strongest as
follows:
\[
\textit{soviel ich wei\ss}
\;<\;
\textit{ja}
\;<\;
\textit{bekanntlich}.
\]

\subsection{Controversy decomposition}

Let
$pc_{\text{prop}} \in [0,1]$ denote propositional controversy and
$pc_{\text{prag}} \in [0,1]$ pragmatic controversy. Overall perceived
controversy is defined as:
\begin{equation}
pc \,=\, PC(pc_{\text{prop}}, pc_{\text{prag}})
\,=\,
pc_{\text{prop}} \cdot pc_{\text{prag}}.
\end{equation}

This multiplicative form captures the idea that a proposition can be globally
contentious while being relatively uncontroversial in the speaker's local
community, and vice versa.

\subsection{Base utility}

Let $g \in [0,1]$ denote speaker goal strength. The speaker's base context
utility is:
\begin{equation}
U(pc,g)
\,=\,
- w_{pc}\,pc
+ w_g\,g
- w_{\mathrm{int}}\,pc\,g,
\label{eq:model-utility}
\end{equation}
where $w_{pc}, w_g, w_{\mathrm{int}} > 0$.

The first term penalizes controversy, the second rewards persuasive effort,
and the interaction term captures that pursuing a strong goal is increasingly
costly when controversy is high.

\subsection{Noisy threshold intervals}

Each marker $u_i$ is associated with a lower threshold $\theta_i$ and an upper
threshold $\theta_{i+1}$, where the upper threshold is the lower threshold of
the next stronger marker. For the strongest marker, the upper boundary is
$+\infty$.

Thresholds are treated as noisy rather than fixed:
\begin{equation}
\theta_i \sim \mathcal{N}(\mu_i, \sigma_i^2),
\end{equation}
with ordered means
\[
\mu_{\textit{soviel ich wei\ss}}
\;<\;
\mu_{\textit{ja}}
\;<\;
\mu_{\textit{bekanntlich}}.
\]

A marker fits the context to the extent that the utility falls inside its
licensed interval. The corresponding noisy interval-fit term is:
\begin{equation}
\mathrm{Fit}_{\mathrm{interval}}(u_i \mid U)
\;\propto\;
\mathbb{E}_{\theta_i}\left[\sigma\!\left(\lambda(U-\theta_i)\right)\right]
\cdot
\mathbb{E}_{\theta_{i+1}}\left[\sigma\!\left(-\lambda(U-\theta_{i+1})\right)\right],
\label{eq:model-fit}
\end{equation}
where $\sigma(x)=(1+e^{-x})^{-1}$ is the logistic function and $\lambda > 0$
controls boundary sharpness.

The first factor rewards utilities above the lower threshold; the second
factor penalizes utilities above the next stronger marker's lower threshold.
Thus, each marker is associated with an interval rather than a single cutoff.

\subsection{Marker costs and sufficiency blocking}

In addition to interval fit, each marker carries a marker-specific pragmatic
cost $c_u \geq 0$. The weakest marker can be treated as the reference level,
with stronger markers incurring larger costs.

Speaker choice is therefore modeled as:
\begin{equation}
P(u \mid pc,g)
\;\propto\;
\mathrm{Fit}_{\mathrm{interval}}\!\left(u \mid U(pc,g)\right)
\cdot e^{-c_u}.
\label{eq:model-choice}
\end{equation}

This cost term introduces a \emph{weakest-sufficient} bias. A stronger marker
may be licensed by the context, but still lose to a weaker marker if the weaker
option already serves the speaker's purpose sufficiently well. This yields the
prediction of \textbf{sufficiency blocking}: in low-controversy contexts, even
under relatively high goal strength, \textit{bekanntlich} can be blocked if
\textit{ja} or even \textit{soviel ich wei\ss} already provides enough force
at lower cost.

\subsection{Listener inversion}

Given an observed marker $u$, the listener inverts the speaker model to infer
latent context variables:
\begin{equation}
P(pc_{\text{prop}}, pc_{\text{prag}}, g \mid u)
\;\propto\;
P\!\left(u \mid PC(pc_{\text{prop}}, pc_{\text{prag}}), g\right)
\,P(pc_{\text{prop}})\,P(pc_{\text{prag}})\,P(g).
\label{eq:model-posterior}
\end{equation}

Because stronger markers generally require higher utility, they shift posterior
mass toward lower controversy and higher goal strength. However, the cost term
weakens a purely monotonic mapping from marker strength to speaker goal:
observing a weaker marker can reflect either lower utility or the fact that a
stronger marker would have been pragmatically excessive.

\subsection{Adoption model}

The listener's downstream adoption decision depends on the posterior over goal
strength and controversy:
\begin{equation}
\text{logit}\left(P_{\mathrm{adopt}}\right)
\;=
\eta_0
+ \eta_g\,\mathbb{E}[g \mid u,pc]
- \eta_{pc}\,pc
- \eta_{\mathrm{int}}\,\mathbb{E}[g \mid u,pc]\,pc.
\label{eq:model-adopt}
\end{equation}

Here, $\eta_0$ is a baseline adoption tendency, $\eta_g$ captures the positive
effect of inferred speaker goal, $\eta_{pc}$ penalizes controversy, and
$\eta_{\mathrm{int}}$ allows controversy to attenuate the effect of inferred
goal strength.

\subsection{Qualitative predictions}

The model yields the following core predictions:
\begin{enumerate}
  \item \textbf{Goal effect:} holding controversy fixed, higher $g$ increases
        the probability of stronger markers.
  \item \textbf{Controversy effect:} holding $g$ fixed, higher $pc$ increases
        the probability of weaker markers.
  \item \textbf{Attenuated goal effect under controversy:} the impact of $g$
        on marker strength is reduced when $pc$ is high.
  \item \textbf{Sufficiency blocking:} in low-$pc$ regions, a stronger marker
        can be blocked by a weaker marker if the weaker one already does enough.
  \item \textbf{Listener-side inversion:} stronger markers tend to induce lower
        inferred controversy and higher inferred goal strength.
  \item \textbf{Credibility discounting:} even when a strong marker supports
        high inferred goal strength, its effect on action adoption is reduced
        in highly controversial contexts.
\end{enumerate}

\subsection{Interpretive summary}

Relative to a vanilla RSA model that directly softmaxes over utterance utility,
this model treats markers as corresponding to \emph{intervals} on a latent
utility scale. The interval-fit component answers the question:
\enquote{How well does this marker match the current utility region?}
The cost term answers the question:
\enquote{Is this marker stronger or more marked than the speaker needs?}
Only their combination determines actual choice.

\end{document}
.
%
% Compile stand-alone:
%   pdflatex model.tex
%
% Embed in main.tex:
%   % model.tex
% Stand-alone description of the cost-augmented noisy-threshold RSA model.
% Can be embedded in main.tex via \input{model}.
%
% Compile stand-alone:
%   pdflatex model.tex
%
% Embed in main.tex:
%   \input{model}

\documentclass[12pt]{article}
\usepackage{CSPArticleStyle}
\usepackage{csquotes}
\usepackage[margin=1in]{geometry}
\usepackage{booktabs}
\usepackage{parskip}
\usepackage{amsmath}
\usepackage{amssymb}
\usepackage{times}
\usepackage[T1]{fontenc}
\usepackage[utf8]{inputenc}
\usepackage{xcolor}

\title{Cost-Augmented Noisy-Threshold RSA Model for Consensus Markers}
\author{}
\date{}

\begin{document}

%% ============================================================
%% When embedding in main.tex, delete from \documentclass to
%% \begin{document} and delete \end{document} at the bottom.
%% ============================================================

\section{Model}

\subsection{Overview}

We model German consensus markers as lexical items that are licensed over
intervals of a latent utility scale. The model combines three ideas:
\begin{enumerate}
  \item perceived controversy is jointly determined by a propositional and a pragmatic component;
  \item speaker choice depends on a utility trade-off between controversy and persuasive goal strength;
  \item stronger markers incur additional pragmatic costs, which creates a
        \emph{weakest-sufficient} bias and can block overly strong markers.
\end{enumerate}

The three investigated markers are ordered from weakest to strongest as
follows:
\[
\textit{soviel ich wei\ss}
\;<\;
\textit{ja}
\;<\;
\textit{bekanntlich}.
\]

\subsection{Controversy decomposition}

Let
$pc_{\text{prop}} \in [0,1]$ denote propositional controversy and
$pc_{\text{prag}} \in [0,1]$ pragmatic controversy. Overall perceived
controversy is defined as:
\begin{equation}
pc \,=\, PC(pc_{\text{prop}}, pc_{\text{prag}})
\,=\,
pc_{\text{prop}} \cdot pc_{\text{prag}}.
\end{equation}

This multiplicative form captures the idea that a proposition can be globally
contentious while being relatively uncontroversial in the speaker's local
community, and vice versa.

\subsection{Base utility}

Let $g \in [0,1]$ denote speaker goal strength. The speaker's base context
utility is:
\begin{equation}
U(pc,g)
\,=\,
- w_{pc}\,pc
+ w_g\,g
- w_{\mathrm{int}}\,pc\,g,
\label{eq:model-utility}
\end{equation}
where $w_{pc}, w_g, w_{\mathrm{int}} > 0$.

The first term penalizes controversy, the second rewards persuasive effort,
and the interaction term captures that pursuing a strong goal is increasingly
costly when controversy is high.

\subsection{Noisy threshold intervals}

Each marker $u_i$ is associated with a lower threshold $\theta_i$ and an upper
threshold $\theta_{i+1}$, where the upper threshold is the lower threshold of
the next stronger marker. For the strongest marker, the upper boundary is
$+\infty$.

Thresholds are treated as noisy rather than fixed:
\begin{equation}
\theta_i \sim \mathcal{N}(\mu_i, \sigma_i^2),
\end{equation}
with ordered means
\[
\mu_{\textit{soviel ich wei\ss}}
\;<\;
\mu_{\textit{ja}}
\;<\;
\mu_{\textit{bekanntlich}}.
\]

A marker fits the context to the extent that the utility falls inside its
licensed interval. The corresponding noisy interval-fit term is:
\begin{equation}
\mathrm{Fit}_{\mathrm{interval}}(u_i \mid U)
\;\propto\;
\mathbb{E}_{\theta_i}\left[\sigma\!\left(\lambda(U-\theta_i)\right)\right]
\cdot
\mathbb{E}_{\theta_{i+1}}\left[\sigma\!\left(-\lambda(U-\theta_{i+1})\right)\right],
\label{eq:model-fit}
\end{equation}
where $\sigma(x)=(1+e^{-x})^{-1}$ is the logistic function and $\lambda > 0$
controls boundary sharpness.

The first factor rewards utilities above the lower threshold; the second
factor penalizes utilities above the next stronger marker's lower threshold.
Thus, each marker is associated with an interval rather than a single cutoff.

\subsection{Marker costs and sufficiency blocking}

In addition to interval fit, each marker carries a marker-specific pragmatic
cost $c_u \geq 0$. The weakest marker can be treated as the reference level,
with stronger markers incurring larger costs.

Speaker choice is therefore modeled as:
\begin{equation}
P(u \mid pc,g)
\;\propto\;
\mathrm{Fit}_{\mathrm{interval}}\!\left(u \mid U(pc,g)\right)
\cdot e^{-c_u}.
\label{eq:model-choice}
\end{equation}

This cost term introduces a \emph{weakest-sufficient} bias. A stronger marker
may be licensed by the context, but still lose to a weaker marker if the weaker
option already serves the speaker's purpose sufficiently well. This yields the
prediction of \textbf{sufficiency blocking}: in low-controversy contexts, even
under relatively high goal strength, \textit{bekanntlich} can be blocked if
\textit{ja} or even \textit{soviel ich wei\ss} already provides enough force
at lower cost.

\subsection{Listener inversion}

Given an observed marker $u$, the listener inverts the speaker model to infer
latent context variables:
\begin{equation}
P(pc_{\text{prop}}, pc_{\text{prag}}, g \mid u)
\;\propto\;
P\!\left(u \mid PC(pc_{\text{prop}}, pc_{\text{prag}}), g\right)
\,P(pc_{\text{prop}})\,P(pc_{\text{prag}})\,P(g).
\label{eq:model-posterior}
\end{equation}

Because stronger markers generally require higher utility, they shift posterior
mass toward lower controversy and higher goal strength. However, the cost term
weakens a purely monotonic mapping from marker strength to speaker goal:
observing a weaker marker can reflect either lower utility or the fact that a
stronger marker would have been pragmatically excessive.

\subsection{Adoption model}

The listener's downstream adoption decision depends on the posterior over goal
strength and controversy:
\begin{equation}
\text{logit}\left(P_{\mathrm{adopt}}\right)
\;=
\eta_0
+ \eta_g\,\mathbb{E}[g \mid u,pc]
- \eta_{pc}\,pc
- \eta_{\mathrm{int}}\,\mathbb{E}[g \mid u,pc]\,pc.
\label{eq:model-adopt}
\end{equation}

Here, $\eta_0$ is a baseline adoption tendency, $\eta_g$ captures the positive
effect of inferred speaker goal, $\eta_{pc}$ penalizes controversy, and
$\eta_{\mathrm{int}}$ allows controversy to attenuate the effect of inferred
goal strength.

\subsection{Qualitative predictions}

The model yields the following core predictions:
\begin{enumerate}
  \item \textbf{Goal effect:} holding controversy fixed, higher $g$ increases
        the probability of stronger markers.
  \item \textbf{Controversy effect:} holding $g$ fixed, higher $pc$ increases
        the probability of weaker markers.
  \item \textbf{Attenuated goal effect under controversy:} the impact of $g$
        on marker strength is reduced when $pc$ is high.
  \item \textbf{Sufficiency blocking:} in low-$pc$ regions, a stronger marker
        can be blocked by a weaker marker if the weaker one already does enough.
  \item \textbf{Listener-side inversion:} stronger markers tend to induce lower
        inferred controversy and higher inferred goal strength.
  \item \textbf{Credibility discounting:} even when a strong marker supports
        high inferred goal strength, its effect on action adoption is reduced
        in highly controversial contexts.
\end{enumerate}

\subsection{Interpretive summary}

Relative to a vanilla RSA model that directly softmaxes over utterance utility,
this model treats markers as corresponding to \emph{intervals} on a latent
utility scale. The interval-fit component answers the question:
\enquote{How well does this marker match the current utility region?}
The cost term answers the question:
\enquote{Is this marker stronger or more marked than the speaker needs?}
Only their combination determines actual choice.

\end{document}


\documentclass[12pt]{article}
\usepackage{CSPArticleStyle}
\usepackage{csquotes}
\usepackage[margin=1in]{geometry}
\usepackage{booktabs}
\usepackage{parskip}
\usepackage{amsmath}
\usepackage{amssymb}
\usepackage{times}
\usepackage[T1]{fontenc}
\usepackage[utf8]{inputenc}
\usepackage{xcolor}

\title{Cost-Augmented Noisy-Threshold RSA Model for Consensus Markers}
\author{}
\date{}

\begin{document}

%% ============================================================
%% When embedding in main.tex, delete from \documentclass to
%% \begin{document} and delete \end{document} at the bottom.
%% ============================================================

\section{Model}

\subsection{Overview}

We model German consensus markers as lexical items that are licensed over
intervals of a latent utility scale. The model combines three ideas:
\begin{enumerate}
  \item perceived controversy is jointly determined by a propositional and a pragmatic component;
  \item speaker choice depends on a utility trade-off between controversy and persuasive goal strength;
  \item stronger markers incur additional pragmatic costs, which creates a
        \emph{weakest-sufficient} bias and can block overly strong markers.
\end{enumerate}

The three investigated markers are ordered from weakest to strongest as
follows:
\[
\textit{soviel ich wei\ss}
\;<\;
\textit{ja}
\;<\;
\textit{bekanntlich}.
\]

\subsection{Controversy decomposition}

Let
$pc_{\text{prop}} \in [0,1]$ denote propositional controversy and
$pc_{\text{prag}} \in [0,1]$ pragmatic controversy. Overall perceived
controversy is defined as:
\begin{equation}
pc \,=\, PC(pc_{\text{prop}}, pc_{\text{prag}})
\,=\,
pc_{\text{prop}} \cdot pc_{\text{prag}}.
\end{equation}

This multiplicative form captures the idea that a proposition can be globally
contentious while being relatively uncontroversial in the speaker's local
community, and vice versa.

\subsection{Base utility}

Let $g \in [0,1]$ denote speaker goal strength. The speaker's base context
utility is:
\begin{equation}
U(pc,g)
\,=\,
- w_{pc}\,pc
+ w_g\,g
- w_{\mathrm{int}}\,pc\,g,
\label{eq:model-utility}
\end{equation}
where $w_{pc}, w_g, w_{\mathrm{int}} > 0$.

The first term penalizes controversy, the second rewards persuasive effort,
and the interaction term captures that pursuing a strong goal is increasingly
costly when controversy is high.

\subsection{Noisy threshold intervals}

Each marker $u_i$ is associated with a lower threshold $\theta_i$ and an upper
threshold $\theta_{i+1}$, where the upper threshold is the lower threshold of
the next stronger marker. For the strongest marker, the upper boundary is
$+\infty$.

Thresholds are treated as noisy rather than fixed:
\begin{equation}
\theta_i \sim \mathcal{N}(\mu_i, \sigma_i^2),
\end{equation}
with ordered means
\[
\mu_{\textit{soviel ich wei\ss}}
\;<\;
\mu_{\textit{ja}}
\;<\;
\mu_{\textit{bekanntlich}}.
\]

A marker fits the context to the extent that the utility falls inside its
licensed interval. The corresponding noisy interval-fit term is:
\begin{equation}
\mathrm{Fit}_{\mathrm{interval}}(u_i \mid U)
\;\propto\;
\mathbb{E}_{\theta_i}\left[\sigma\!\left(\lambda(U-\theta_i)\right)\right]
\cdot
\mathbb{E}_{\theta_{i+1}}\left[\sigma\!\left(-\lambda(U-\theta_{i+1})\right)\right],
\label{eq:model-fit}
\end{equation}
where $\sigma(x)=(1+e^{-x})^{-1}$ is the logistic function and $\lambda > 0$
controls boundary sharpness.

The first factor rewards utilities above the lower threshold; the second
factor penalizes utilities above the next stronger marker's lower threshold.
Thus, each marker is associated with an interval rather than a single cutoff.

\subsection{Marker costs and sufficiency blocking}

In addition to interval fit, each marker carries a marker-specific pragmatic
cost $c_u \geq 0$. The weakest marker can be treated as the reference level,
with stronger markers incurring larger costs.

Speaker choice is therefore modeled as:
\begin{equation}
P(u \mid pc,g)
\;\propto\;
\mathrm{Fit}_{\mathrm{interval}}\!\left(u \mid U(pc,g)\right)
\cdot e^{-c_u}.
\label{eq:model-choice}
\end{equation}

This cost term introduces a \emph{weakest-sufficient} bias. A stronger marker
may be licensed by the context, but still lose to a weaker marker if the weaker
option already serves the speaker's purpose sufficiently well. This yields the
prediction of \textbf{sufficiency blocking}: in low-controversy contexts, even
under relatively high goal strength, \textit{bekanntlich} can be blocked if
\textit{ja} or even \textit{soviel ich wei\ss} already provides enough force
at lower cost.

\subsection{Listener inversion}

Given an observed marker $u$, the listener inverts the speaker model to infer
latent context variables:
\begin{equation}
P(pc_{\text{prop}}, pc_{\text{prag}}, g \mid u)
\;\propto\;
P\!\left(u \mid PC(pc_{\text{prop}}, pc_{\text{prag}}), g\right)
\,P(pc_{\text{prop}})\,P(pc_{\text{prag}})\,P(g).
\label{eq:model-posterior}
\end{equation}

Because stronger markers generally require higher utility, they shift posterior
mass toward lower controversy and higher goal strength. However, the cost term
weakens a purely monotonic mapping from marker strength to speaker goal:
observing a weaker marker can reflect either lower utility or the fact that a
stronger marker would have been pragmatically excessive.

\subsection{Adoption model}

The listener's downstream adoption decision depends on the posterior over goal
strength and controversy:
\begin{equation}
\text{logit}\left(P_{\mathrm{adopt}}\right)
\;=
\eta_0
+ \eta_g\,\mathbb{E}[g \mid u,pc]
- \eta_{pc}\,pc
- \eta_{\mathrm{int}}\,\mathbb{E}[g \mid u,pc]\,pc.
\label{eq:model-adopt}
\end{equation}

Here, $\eta_0$ is a baseline adoption tendency, $\eta_g$ captures the positive
effect of inferred speaker goal, $\eta_{pc}$ penalizes controversy, and
$\eta_{\mathrm{int}}$ allows controversy to attenuate the effect of inferred
goal strength.

\subsection{Qualitative predictions}

The model yields the following core predictions:
\begin{enumerate}
  \item \textbf{Goal effect:} holding controversy fixed, higher $g$ increases
        the probability of stronger markers.
  \item \textbf{Controversy effect:} holding $g$ fixed, higher $pc$ increases
        the probability of weaker markers.
  \item \textbf{Attenuated goal effect under controversy:} the impact of $g$
        on marker strength is reduced when $pc$ is high.
  \item \textbf{Sufficiency blocking:} in low-$pc$ regions, a stronger marker
        can be blocked by a weaker marker if the weaker one already does enough.
  \item \textbf{Listener-side inversion:} stronger markers tend to induce lower
        inferred controversy and higher inferred goal strength.
  \item \textbf{Credibility discounting:} even when a strong marker supports
        high inferred goal strength, its effect on action adoption is reduced
        in highly controversial contexts.
\end{enumerate}

\subsection{Interpretive summary}

Relative to a vanilla RSA model that directly softmaxes over utterance utility,
this model treats markers as corresponding to \emph{intervals} on a latent
utility scale. The interval-fit component answers the question:
\enquote{How well does this marker match the current utility region?}
The cost term answers the question:
\enquote{Is this marker stronger or more marked than the speaker needs?}
Only their combination determines actual choice.

\end{document}
.
%
% Compile stand-alone:
%   pdflatex model.tex
%
% Embed in main.tex:
%   % model.tex
% Stand-alone description of the cost-augmented noisy-threshold RSA model.
% Can be embedded in main.tex via % model.tex
% Stand-alone description of the cost-augmented noisy-threshold RSA model.
% Can be embedded in main.tex via \input{model}.
%
% Compile stand-alone:
%   pdflatex model.tex
%
% Embed in main.tex:
%   \input{model}

\documentclass[12pt]{article}
\usepackage{CSPArticleStyle}
\usepackage{csquotes}
\usepackage[margin=1in]{geometry}
\usepackage{booktabs}
\usepackage{parskip}
\usepackage{amsmath}
\usepackage{amssymb}
\usepackage{times}
\usepackage[T1]{fontenc}
\usepackage[utf8]{inputenc}
\usepackage{xcolor}

\title{Cost-Augmented Noisy-Threshold RSA Model for Consensus Markers}
\author{}
\date{}

\begin{document}

%% ============================================================
%% When embedding in main.tex, delete from \documentclass to
%% \begin{document} and delete \end{document} at the bottom.
%% ============================================================

\section{Model}

\subsection{Overview}

We model German consensus markers as lexical items that are licensed over
intervals of a latent utility scale. The model combines three ideas:
\begin{enumerate}
  \item perceived controversy is jointly determined by a propositional and a pragmatic component;
  \item speaker choice depends on a utility trade-off between controversy and persuasive goal strength;
  \item stronger markers incur additional pragmatic costs, which creates a
        \emph{weakest-sufficient} bias and can block overly strong markers.
\end{enumerate}

The three investigated markers are ordered from weakest to strongest as
follows:
\[
\textit{soviel ich wei\ss}
\;<\;
\textit{ja}
\;<\;
\textit{bekanntlich}.
\]

\subsection{Controversy decomposition}

Let
$pc_{\text{prop}} \in [0,1]$ denote propositional controversy and
$pc_{\text{prag}} \in [0,1]$ pragmatic controversy. Overall perceived
controversy is defined as:
\begin{equation}
pc \,=\, PC(pc_{\text{prop}}, pc_{\text{prag}})
\,=\,
pc_{\text{prop}} \cdot pc_{\text{prag}}.
\end{equation}

This multiplicative form captures the idea that a proposition can be globally
contentious while being relatively uncontroversial in the speaker's local
community, and vice versa.

\subsection{Base utility}

Let $g \in [0,1]$ denote speaker goal strength. The speaker's base context
utility is:
\begin{equation}
U(pc,g)
\,=\,
- w_{pc}\,pc
+ w_g\,g
- w_{\mathrm{int}}\,pc\,g,
\label{eq:model-utility}
\end{equation}
where $w_{pc}, w_g, w_{\mathrm{int}} > 0$.

The first term penalizes controversy, the second rewards persuasive effort,
and the interaction term captures that pursuing a strong goal is increasingly
costly when controversy is high.

\subsection{Noisy threshold intervals}

Each marker $u_i$ is associated with a lower threshold $\theta_i$ and an upper
threshold $\theta_{i+1}$, where the upper threshold is the lower threshold of
the next stronger marker. For the strongest marker, the upper boundary is
$+\infty$.

Thresholds are treated as noisy rather than fixed:
\begin{equation}
\theta_i \sim \mathcal{N}(\mu_i, \sigma_i^2),
\end{equation}
with ordered means
\[
\mu_{\textit{soviel ich wei\ss}}
\;<\;
\mu_{\textit{ja}}
\;<\;
\mu_{\textit{bekanntlich}}.
\]

A marker fits the context to the extent that the utility falls inside its
licensed interval. The corresponding noisy interval-fit term is:
\begin{equation}
\mathrm{Fit}_{\mathrm{interval}}(u_i \mid U)
\;\propto\;
\mathbb{E}_{\theta_i}\left[\sigma\!\left(\lambda(U-\theta_i)\right)\right]
\cdot
\mathbb{E}_{\theta_{i+1}}\left[\sigma\!\left(-\lambda(U-\theta_{i+1})\right)\right],
\label{eq:model-fit}
\end{equation}
where $\sigma(x)=(1+e^{-x})^{-1}$ is the logistic function and $\lambda > 0$
controls boundary sharpness.

The first factor rewards utilities above the lower threshold; the second
factor penalizes utilities above the next stronger marker's lower threshold.
Thus, each marker is associated with an interval rather than a single cutoff.

\subsection{Marker costs and sufficiency blocking}

In addition to interval fit, each marker carries a marker-specific pragmatic
cost $c_u \geq 0$. The weakest marker can be treated as the reference level,
with stronger markers incurring larger costs.

Speaker choice is therefore modeled as:
\begin{equation}
P(u \mid pc,g)
\;\propto\;
\mathrm{Fit}_{\mathrm{interval}}\!\left(u \mid U(pc,g)\right)
\cdot e^{-c_u}.
\label{eq:model-choice}
\end{equation}

This cost term introduces a \emph{weakest-sufficient} bias. A stronger marker
may be licensed by the context, but still lose to a weaker marker if the weaker
option already serves the speaker's purpose sufficiently well. This yields the
prediction of \textbf{sufficiency blocking}: in low-controversy contexts, even
under relatively high goal strength, \textit{bekanntlich} can be blocked if
\textit{ja} or even \textit{soviel ich wei\ss} already provides enough force
at lower cost.

\subsection{Listener inversion}

Given an observed marker $u$, the listener inverts the speaker model to infer
latent context variables:
\begin{equation}
P(pc_{\text{prop}}, pc_{\text{prag}}, g \mid u)
\;\propto\;
P\!\left(u \mid PC(pc_{\text{prop}}, pc_{\text{prag}}), g\right)
\,P(pc_{\text{prop}})\,P(pc_{\text{prag}})\,P(g).
\label{eq:model-posterior}
\end{equation}

Because stronger markers generally require higher utility, they shift posterior
mass toward lower controversy and higher goal strength. However, the cost term
weakens a purely monotonic mapping from marker strength to speaker goal:
observing a weaker marker can reflect either lower utility or the fact that a
stronger marker would have been pragmatically excessive.

\subsection{Adoption model}

The listener's downstream adoption decision depends on the posterior over goal
strength and controversy:
\begin{equation}
\text{logit}\left(P_{\mathrm{adopt}}\right)
\;=
\eta_0
+ \eta_g\,\mathbb{E}[g \mid u,pc]
- \eta_{pc}\,pc
- \eta_{\mathrm{int}}\,\mathbb{E}[g \mid u,pc]\,pc.
\label{eq:model-adopt}
\end{equation}

Here, $\eta_0$ is a baseline adoption tendency, $\eta_g$ captures the positive
effect of inferred speaker goal, $\eta_{pc}$ penalizes controversy, and
$\eta_{\mathrm{int}}$ allows controversy to attenuate the effect of inferred
goal strength.

\subsection{Qualitative predictions}

The model yields the following core predictions:
\begin{enumerate}
  \item \textbf{Goal effect:} holding controversy fixed, higher $g$ increases
        the probability of stronger markers.
  \item \textbf{Controversy effect:} holding $g$ fixed, higher $pc$ increases
        the probability of weaker markers.
  \item \textbf{Attenuated goal effect under controversy:} the impact of $g$
        on marker strength is reduced when $pc$ is high.
  \item \textbf{Sufficiency blocking:} in low-$pc$ regions, a stronger marker
        can be blocked by a weaker marker if the weaker one already does enough.
  \item \textbf{Listener-side inversion:} stronger markers tend to induce lower
        inferred controversy and higher inferred goal strength.
  \item \textbf{Credibility discounting:} even when a strong marker supports
        high inferred goal strength, its effect on action adoption is reduced
        in highly controversial contexts.
\end{enumerate}

\subsection{Interpretive summary}

Relative to a vanilla RSA model that directly softmaxes over utterance utility,
this model treats markers as corresponding to \emph{intervals} on a latent
utility scale. The interval-fit component answers the question:
\enquote{How well does this marker match the current utility region?}
The cost term answers the question:
\enquote{Is this marker stronger or more marked than the speaker needs?}
Only their combination determines actual choice.

\end{document}
.
%
% Compile stand-alone:
%   pdflatex model.tex
%
% Embed in main.tex:
%   % model.tex
% Stand-alone description of the cost-augmented noisy-threshold RSA model.
% Can be embedded in main.tex via \input{model}.
%
% Compile stand-alone:
%   pdflatex model.tex
%
% Embed in main.tex:
%   \input{model}

\documentclass[12pt]{article}
\usepackage{CSPArticleStyle}
\usepackage{csquotes}
\usepackage[margin=1in]{geometry}
\usepackage{booktabs}
\usepackage{parskip}
\usepackage{amsmath}
\usepackage{amssymb}
\usepackage{times}
\usepackage[T1]{fontenc}
\usepackage[utf8]{inputenc}
\usepackage{xcolor}

\title{Cost-Augmented Noisy-Threshold RSA Model for Consensus Markers}
\author{}
\date{}

\begin{document}

%% ============================================================
%% When embedding in main.tex, delete from \documentclass to
%% \begin{document} and delete \end{document} at the bottom.
%% ============================================================

\section{Model}

\subsection{Overview}

We model German consensus markers as lexical items that are licensed over
intervals of a latent utility scale. The model combines three ideas:
\begin{enumerate}
  \item perceived controversy is jointly determined by a propositional and a pragmatic component;
  \item speaker choice depends on a utility trade-off between controversy and persuasive goal strength;
  \item stronger markers incur additional pragmatic costs, which creates a
        \emph{weakest-sufficient} bias and can block overly strong markers.
\end{enumerate}

The three investigated markers are ordered from weakest to strongest as
follows:
\[
\textit{soviel ich wei\ss}
\;<\;
\textit{ja}
\;<\;
\textit{bekanntlich}.
\]

\subsection{Controversy decomposition}

Let
$pc_{\text{prop}} \in [0,1]$ denote propositional controversy and
$pc_{\text{prag}} \in [0,1]$ pragmatic controversy. Overall perceived
controversy is defined as:
\begin{equation}
pc \,=\, PC(pc_{\text{prop}}, pc_{\text{prag}})
\,=\,
pc_{\text{prop}} \cdot pc_{\text{prag}}.
\end{equation}

This multiplicative form captures the idea that a proposition can be globally
contentious while being relatively uncontroversial in the speaker's local
community, and vice versa.

\subsection{Base utility}

Let $g \in [0,1]$ denote speaker goal strength. The speaker's base context
utility is:
\begin{equation}
U(pc,g)
\,=\,
- w_{pc}\,pc
+ w_g\,g
- w_{\mathrm{int}}\,pc\,g,
\label{eq:model-utility}
\end{equation}
where $w_{pc}, w_g, w_{\mathrm{int}} > 0$.

The first term penalizes controversy, the second rewards persuasive effort,
and the interaction term captures that pursuing a strong goal is increasingly
costly when controversy is high.

\subsection{Noisy threshold intervals}

Each marker $u_i$ is associated with a lower threshold $\theta_i$ and an upper
threshold $\theta_{i+1}$, where the upper threshold is the lower threshold of
the next stronger marker. For the strongest marker, the upper boundary is
$+\infty$.

Thresholds are treated as noisy rather than fixed:
\begin{equation}
\theta_i \sim \mathcal{N}(\mu_i, \sigma_i^2),
\end{equation}
with ordered means
\[
\mu_{\textit{soviel ich wei\ss}}
\;<\;
\mu_{\textit{ja}}
\;<\;
\mu_{\textit{bekanntlich}}.
\]

A marker fits the context to the extent that the utility falls inside its
licensed interval. The corresponding noisy interval-fit term is:
\begin{equation}
\mathrm{Fit}_{\mathrm{interval}}(u_i \mid U)
\;\propto\;
\mathbb{E}_{\theta_i}\left[\sigma\!\left(\lambda(U-\theta_i)\right)\right]
\cdot
\mathbb{E}_{\theta_{i+1}}\left[\sigma\!\left(-\lambda(U-\theta_{i+1})\right)\right],
\label{eq:model-fit}
\end{equation}
where $\sigma(x)=(1+e^{-x})^{-1}$ is the logistic function and $\lambda > 0$
controls boundary sharpness.

The first factor rewards utilities above the lower threshold; the second
factor penalizes utilities above the next stronger marker's lower threshold.
Thus, each marker is associated with an interval rather than a single cutoff.

\subsection{Marker costs and sufficiency blocking}

In addition to interval fit, each marker carries a marker-specific pragmatic
cost $c_u \geq 0$. The weakest marker can be treated as the reference level,
with stronger markers incurring larger costs.

Speaker choice is therefore modeled as:
\begin{equation}
P(u \mid pc,g)
\;\propto\;
\mathrm{Fit}_{\mathrm{interval}}\!\left(u \mid U(pc,g)\right)
\cdot e^{-c_u}.
\label{eq:model-choice}
\end{equation}

This cost term introduces a \emph{weakest-sufficient} bias. A stronger marker
may be licensed by the context, but still lose to a weaker marker if the weaker
option already serves the speaker's purpose sufficiently well. This yields the
prediction of \textbf{sufficiency blocking}: in low-controversy contexts, even
under relatively high goal strength, \textit{bekanntlich} can be blocked if
\textit{ja} or even \textit{soviel ich wei\ss} already provides enough force
at lower cost.

\subsection{Listener inversion}

Given an observed marker $u$, the listener inverts the speaker model to infer
latent context variables:
\begin{equation}
P(pc_{\text{prop}}, pc_{\text{prag}}, g \mid u)
\;\propto\;
P\!\left(u \mid PC(pc_{\text{prop}}, pc_{\text{prag}}), g\right)
\,P(pc_{\text{prop}})\,P(pc_{\text{prag}})\,P(g).
\label{eq:model-posterior}
\end{equation}

Because stronger markers generally require higher utility, they shift posterior
mass toward lower controversy and higher goal strength. However, the cost term
weakens a purely monotonic mapping from marker strength to speaker goal:
observing a weaker marker can reflect either lower utility or the fact that a
stronger marker would have been pragmatically excessive.

\subsection{Adoption model}

The listener's downstream adoption decision depends on the posterior over goal
strength and controversy:
\begin{equation}
\text{logit}\left(P_{\mathrm{adopt}}\right)
\;=
\eta_0
+ \eta_g\,\mathbb{E}[g \mid u,pc]
- \eta_{pc}\,pc
- \eta_{\mathrm{int}}\,\mathbb{E}[g \mid u,pc]\,pc.
\label{eq:model-adopt}
\end{equation}

Here, $\eta_0$ is a baseline adoption tendency, $\eta_g$ captures the positive
effect of inferred speaker goal, $\eta_{pc}$ penalizes controversy, and
$\eta_{\mathrm{int}}$ allows controversy to attenuate the effect of inferred
goal strength.

\subsection{Qualitative predictions}

The model yields the following core predictions:
\begin{enumerate}
  \item \textbf{Goal effect:} holding controversy fixed, higher $g$ increases
        the probability of stronger markers.
  \item \textbf{Controversy effect:} holding $g$ fixed, higher $pc$ increases
        the probability of weaker markers.
  \item \textbf{Attenuated goal effect under controversy:} the impact of $g$
        on marker strength is reduced when $pc$ is high.
  \item \textbf{Sufficiency blocking:} in low-$pc$ regions, a stronger marker
        can be blocked by a weaker marker if the weaker one already does enough.
  \item \textbf{Listener-side inversion:} stronger markers tend to induce lower
        inferred controversy and higher inferred goal strength.
  \item \textbf{Credibility discounting:} even when a strong marker supports
        high inferred goal strength, its effect on action adoption is reduced
        in highly controversial contexts.
\end{enumerate}

\subsection{Interpretive summary}

Relative to a vanilla RSA model that directly softmaxes over utterance utility,
this model treats markers as corresponding to \emph{intervals} on a latent
utility scale. The interval-fit component answers the question:
\enquote{How well does this marker match the current utility region?}
The cost term answers the question:
\enquote{Is this marker stronger or more marked than the speaker needs?}
Only their combination determines actual choice.

\end{document}


\documentclass[12pt]{article}
\usepackage{CSPArticleStyle}
\usepackage{csquotes}
\usepackage[margin=1in]{geometry}
\usepackage{booktabs}
\usepackage{parskip}
\usepackage{amsmath}
\usepackage{amssymb}
\usepackage{times}
\usepackage[T1]{fontenc}
\usepackage[utf8]{inputenc}
\usepackage{xcolor}

\title{Cost-Augmented Noisy-Threshold RSA Model for Consensus Markers}
\author{}
\date{}

\begin{document}

%% ============================================================
%% When embedding in main.tex, delete from \documentclass to
%% \begin{document} and delete \end{document} at the bottom.
%% ============================================================

\section{Model}

\subsection{Overview}

We model German consensus markers as lexical items that are licensed over
intervals of a latent utility scale. The model combines three ideas:
\begin{enumerate}
  \item perceived controversy is jointly determined by a propositional and a pragmatic component;
  \item speaker choice depends on a utility trade-off between controversy and persuasive goal strength;
  \item stronger markers incur additional pragmatic costs, which creates a
        \emph{weakest-sufficient} bias and can block overly strong markers.
\end{enumerate}

The three investigated markers are ordered from weakest to strongest as
follows:
\[
\textit{soviel ich wei\ss}
\;<\;
\textit{ja}
\;<\;
\textit{bekanntlich}.
\]

\subsection{Controversy decomposition}

Let
$pc_{\text{prop}} \in [0,1]$ denote propositional controversy and
$pc_{\text{prag}} \in [0,1]$ pragmatic controversy. Overall perceived
controversy is defined as:
\begin{equation}
pc \,=\, PC(pc_{\text{prop}}, pc_{\text{prag}})
\,=\,
pc_{\text{prop}} \cdot pc_{\text{prag}}.
\end{equation}

This multiplicative form captures the idea that a proposition can be globally
contentious while being relatively uncontroversial in the speaker's local
community, and vice versa.

\subsection{Base utility}

Let $g \in [0,1]$ denote speaker goal strength. The speaker's base context
utility is:
\begin{equation}
U(pc,g)
\,=\,
- w_{pc}\,pc
+ w_g\,g
- w_{\mathrm{int}}\,pc\,g,
\label{eq:model-utility}
\end{equation}
where $w_{pc}, w_g, w_{\mathrm{int}} > 0$.

The first term penalizes controversy, the second rewards persuasive effort,
and the interaction term captures that pursuing a strong goal is increasingly
costly when controversy is high.

\subsection{Noisy threshold intervals}

Each marker $u_i$ is associated with a lower threshold $\theta_i$ and an upper
threshold $\theta_{i+1}$, where the upper threshold is the lower threshold of
the next stronger marker. For the strongest marker, the upper boundary is
$+\infty$.

Thresholds are treated as noisy rather than fixed:
\begin{equation}
\theta_i \sim \mathcal{N}(\mu_i, \sigma_i^2),
\end{equation}
with ordered means
\[
\mu_{\textit{soviel ich wei\ss}}
\;<\;
\mu_{\textit{ja}}
\;<\;
\mu_{\textit{bekanntlich}}.
\]

A marker fits the context to the extent that the utility falls inside its
licensed interval. The corresponding noisy interval-fit term is:
\begin{equation}
\mathrm{Fit}_{\mathrm{interval}}(u_i \mid U)
\;\propto\;
\mathbb{E}_{\theta_i}\left[\sigma\!\left(\lambda(U-\theta_i)\right)\right]
\cdot
\mathbb{E}_{\theta_{i+1}}\left[\sigma\!\left(-\lambda(U-\theta_{i+1})\right)\right],
\label{eq:model-fit}
\end{equation}
where $\sigma(x)=(1+e^{-x})^{-1}$ is the logistic function and $\lambda > 0$
controls boundary sharpness.

The first factor rewards utilities above the lower threshold; the second
factor penalizes utilities above the next stronger marker's lower threshold.
Thus, each marker is associated with an interval rather than a single cutoff.

\subsection{Marker costs and sufficiency blocking}

In addition to interval fit, each marker carries a marker-specific pragmatic
cost $c_u \geq 0$. The weakest marker can be treated as the reference level,
with stronger markers incurring larger costs.

Speaker choice is therefore modeled as:
\begin{equation}
P(u \mid pc,g)
\;\propto\;
\mathrm{Fit}_{\mathrm{interval}}\!\left(u \mid U(pc,g)\right)
\cdot e^{-c_u}.
\label{eq:model-choice}
\end{equation}

This cost term introduces a \emph{weakest-sufficient} bias. A stronger marker
may be licensed by the context, but still lose to a weaker marker if the weaker
option already serves the speaker's purpose sufficiently well. This yields the
prediction of \textbf{sufficiency blocking}: in low-controversy contexts, even
under relatively high goal strength, \textit{bekanntlich} can be blocked if
\textit{ja} or even \textit{soviel ich wei\ss} already provides enough force
at lower cost.

\subsection{Listener inversion}

Given an observed marker $u$, the listener inverts the speaker model to infer
latent context variables:
\begin{equation}
P(pc_{\text{prop}}, pc_{\text{prag}}, g \mid u)
\;\propto\;
P\!\left(u \mid PC(pc_{\text{prop}}, pc_{\text{prag}}), g\right)
\,P(pc_{\text{prop}})\,P(pc_{\text{prag}})\,P(g).
\label{eq:model-posterior}
\end{equation}

Because stronger markers generally require higher utility, they shift posterior
mass toward lower controversy and higher goal strength. However, the cost term
weakens a purely monotonic mapping from marker strength to speaker goal:
observing a weaker marker can reflect either lower utility or the fact that a
stronger marker would have been pragmatically excessive.

\subsection{Adoption model}

The listener's downstream adoption decision depends on the posterior over goal
strength and controversy:
\begin{equation}
\text{logit}\left(P_{\mathrm{adopt}}\right)
\;=
\eta_0
+ \eta_g\,\mathbb{E}[g \mid u,pc]
- \eta_{pc}\,pc
- \eta_{\mathrm{int}}\,\mathbb{E}[g \mid u,pc]\,pc.
\label{eq:model-adopt}
\end{equation}

Here, $\eta_0$ is a baseline adoption tendency, $\eta_g$ captures the positive
effect of inferred speaker goal, $\eta_{pc}$ penalizes controversy, and
$\eta_{\mathrm{int}}$ allows controversy to attenuate the effect of inferred
goal strength.

\subsection{Qualitative predictions}

The model yields the following core predictions:
\begin{enumerate}
  \item \textbf{Goal effect:} holding controversy fixed, higher $g$ increases
        the probability of stronger markers.
  \item \textbf{Controversy effect:} holding $g$ fixed, higher $pc$ increases
        the probability of weaker markers.
  \item \textbf{Attenuated goal effect under controversy:} the impact of $g$
        on marker strength is reduced when $pc$ is high.
  \item \textbf{Sufficiency blocking:} in low-$pc$ regions, a stronger marker
        can be blocked by a weaker marker if the weaker one already does enough.
  \item \textbf{Listener-side inversion:} stronger markers tend to induce lower
        inferred controversy and higher inferred goal strength.
  \item \textbf{Credibility discounting:} even when a strong marker supports
        high inferred goal strength, its effect on action adoption is reduced
        in highly controversial contexts.
\end{enumerate}

\subsection{Interpretive summary}

Relative to a vanilla RSA model that directly softmaxes over utterance utility,
this model treats markers as corresponding to \emph{intervals} on a latent
utility scale. The interval-fit component answers the question:
\enquote{How well does this marker match the current utility region?}
The cost term answers the question:
\enquote{Is this marker stronger or more marked than the speaker needs?}
Only their combination determines actual choice.

\end{document}


\documentclass[12pt]{article}
\usepackage{CSPArticleStyle}
\usepackage{csquotes}
\usepackage[margin=1in]{geometry}
\usepackage{booktabs}
\usepackage{parskip}
\usepackage{amsmath}
\usepackage{amssymb}
\usepackage{times}
\usepackage[T1]{fontenc}
\usepackage[utf8]{inputenc}
\usepackage{xcolor}

\title{Cost-Augmented Noisy-Threshold RSA Model for Consensus Markers}
\author{}
\date{}

\begin{document}

%% ============================================================
%% When embedding in main.tex, delete from \documentclass to
%% \begin{document} and delete \end{document} at the bottom.
%% ============================================================

\section{Model}

\subsection{Overview}

We model German consensus markers as lexical items that are licensed over
intervals of a latent utility scale. The model combines three ideas:
\begin{enumerate}
  \item perceived controversy is jointly determined by a propositional and a pragmatic component;
  \item speaker choice depends on a utility trade-off between controversy and persuasive goal strength;
  \item stronger markers incur additional pragmatic costs, which creates a
        \emph{weakest-sufficient} bias and can block overly strong markers.
\end{enumerate}

The three investigated markers are ordered from weakest to strongest as
follows:
\[
\textit{soviel ich wei\ss}
\;<\;
\textit{ja}
\;<\;
\textit{bekanntlich}.
\]

\subsection{Controversy decomposition}

Let
$pc_{\text{prop}} \in [0,1]$ denote propositional controversy and
$pc_{\text{prag}} \in [0,1]$ pragmatic controversy. Overall perceived
controversy is defined as:
\begin{equation}
pc \,=\, PC(pc_{\text{prop}}, pc_{\text{prag}})
\,=\,
pc_{\text{prop}} \cdot pc_{\text{prag}}.
\end{equation}

This multiplicative form captures the idea that a proposition can be globally
contentious while being relatively uncontroversial in the speaker's local
community, and vice versa.

\subsection{Base utility}

Let $g \in [0,1]$ denote speaker goal strength. The speaker's base context
utility is:
\begin{equation}
U(pc,g)
\,=\,
- w_{pc}\,pc
+ w_g\,g
- w_{\mathrm{int}}\,pc\,g,
\label{eq:model-utility}
\end{equation}
where $w_{pc}, w_g, w_{\mathrm{int}} > 0$.

The first term penalizes controversy, the second rewards persuasive effort,
and the interaction term captures that pursuing a strong goal is increasingly
costly when controversy is high.

\subsection{Noisy threshold intervals}

Each marker $u_i$ is associated with a lower threshold $\theta_i$ and an upper
threshold $\theta_{i+1}$, where the upper threshold is the lower threshold of
the next stronger marker. For the strongest marker, the upper boundary is
$+\infty$.

Thresholds are treated as noisy rather than fixed:
\begin{equation}
\theta_i \sim \mathcal{N}(\mu_i, \sigma_i^2),
\end{equation}
with ordered means
\[
\mu_{\textit{soviel ich wei\ss}}
\;<\;
\mu_{\textit{ja}}
\;<\;
\mu_{\textit{bekanntlich}}.
\]

A marker fits the context to the extent that the utility falls inside its
licensed interval. The corresponding noisy interval-fit term is:
\begin{equation}
\mathrm{Fit}_{\mathrm{interval}}(u_i \mid U)
\;\propto\;
\mathbb{E}_{\theta_i}\left[\sigma\!\left(\lambda(U-\theta_i)\right)\right]
\cdot
\mathbb{E}_{\theta_{i+1}}\left[\sigma\!\left(-\lambda(U-\theta_{i+1})\right)\right],
\label{eq:model-fit}
\end{equation}
where $\sigma(x)=(1+e^{-x})^{-1}$ is the logistic function and $\lambda > 0$
controls boundary sharpness.

The first factor rewards utilities above the lower threshold; the second
factor penalizes utilities above the next stronger marker's lower threshold.
Thus, each marker is associated with an interval rather than a single cutoff.

\subsection{Marker costs and sufficiency blocking}

In addition to interval fit, each marker carries a marker-specific pragmatic
cost $c_u \geq 0$. The weakest marker can be treated as the reference level,
with stronger markers incurring larger costs.

Speaker choice is therefore modeled as:
\begin{equation}
P(u \mid pc,g)
\;\propto\;
\mathrm{Fit}_{\mathrm{interval}}\!\left(u \mid U(pc,g)\right)
\cdot e^{-c_u}.
\label{eq:model-choice}
\end{equation}

This cost term introduces a \emph{weakest-sufficient} bias. A stronger marker
may be licensed by the context, but still lose to a weaker marker if the weaker
option already serves the speaker's purpose sufficiently well. This yields the
prediction of \textbf{sufficiency blocking}: in low-controversy contexts, even
under relatively high goal strength, \textit{bekanntlich} can be blocked if
\textit{ja} or even \textit{soviel ich wei\ss} already provides enough force
at lower cost.

\subsection{Listener inversion}

Given an observed marker $u$, the listener inverts the speaker model to infer
latent context variables:
\begin{equation}
P(pc_{\text{prop}}, pc_{\text{prag}}, g \mid u)
\;\propto\;
P\!\left(u \mid PC(pc_{\text{prop}}, pc_{\text{prag}}), g\right)
\,P(pc_{\text{prop}})\,P(pc_{\text{prag}})\,P(g).
\label{eq:model-posterior}
\end{equation}

Because stronger markers generally require higher utility, they shift posterior
mass toward lower controversy and higher goal strength. However, the cost term
weakens a purely monotonic mapping from marker strength to speaker goal:
observing a weaker marker can reflect either lower utility or the fact that a
stronger marker would have been pragmatically excessive.

\subsection{Adoption model}

The listener's downstream adoption decision depends on the posterior over goal
strength and controversy:
\begin{equation}
\text{logit}\left(P_{\mathrm{adopt}}\right)
\;=
\eta_0
+ \eta_g\,\mathbb{E}[g \mid u,pc]
- \eta_{pc}\,pc
- \eta_{\mathrm{int}}\,\mathbb{E}[g \mid u,pc]\,pc.
\label{eq:model-adopt}
\end{equation}

Here, $\eta_0$ is a baseline adoption tendency, $\eta_g$ captures the positive
effect of inferred speaker goal, $\eta_{pc}$ penalizes controversy, and
$\eta_{\mathrm{int}}$ allows controversy to attenuate the effect of inferred
goal strength.

\subsection{Qualitative predictions}

The model yields the following core predictions:
\begin{enumerate}
  \item \textbf{Goal effect:} holding controversy fixed, higher $g$ increases
        the probability of stronger markers.
  \item \textbf{Controversy effect:} holding $g$ fixed, higher $pc$ increases
        the probability of weaker markers.
  \item \textbf{Attenuated goal effect under controversy:} the impact of $g$
        on marker strength is reduced when $pc$ is high.
  \item \textbf{Sufficiency blocking:} in low-$pc$ regions, a stronger marker
        can be blocked by a weaker marker if the weaker one already does enough.
  \item \textbf{Listener-side inversion:} stronger markers tend to induce lower
        inferred controversy and higher inferred goal strength.
  \item \textbf{Credibility discounting:} even when a strong marker supports
        high inferred goal strength, its effect on action adoption is reduced
        in highly controversial contexts.
\end{enumerate}

\subsection{Interpretive summary}

Relative to a vanilla RSA model that directly softmaxes over utterance utility,
this model treats markers as corresponding to \emph{intervals} on a latent
utility scale. The interval-fit component answers the question:
\enquote{How well does this marker match the current utility region?}
The cost term answers the question:
\enquote{Is this marker stronger or more marked than the speaker needs?}
Only their combination determines actual choice.

\end{document}
.
%
% Compile stand-alone:
%   pdflatex model.tex
%
% Embed in main.tex:
%   % model.tex
% Stand-alone description of the cost-augmented noisy-threshold RSA model.
% Can be embedded in main.tex via % model.tex
% Stand-alone description of the cost-augmented noisy-threshold RSA model.
% Can be embedded in main.tex via % model.tex
% Stand-alone description of the cost-augmented noisy-threshold RSA model.
% Can be embedded in main.tex via \input{model}.
%
% Compile stand-alone:
%   pdflatex model.tex
%
% Embed in main.tex:
%   \input{model}

\documentclass[12pt]{article}
\usepackage{CSPArticleStyle}
\usepackage{csquotes}
\usepackage[margin=1in]{geometry}
\usepackage{booktabs}
\usepackage{parskip}
\usepackage{amsmath}
\usepackage{amssymb}
\usepackage{times}
\usepackage[T1]{fontenc}
\usepackage[utf8]{inputenc}
\usepackage{xcolor}

\title{Cost-Augmented Noisy-Threshold RSA Model for Consensus Markers}
\author{}
\date{}

\begin{document}

%% ============================================================
%% When embedding in main.tex, delete from \documentclass to
%% \begin{document} and delete \end{document} at the bottom.
%% ============================================================

\section{Model}

\subsection{Overview}

We model German consensus markers as lexical items that are licensed over
intervals of a latent utility scale. The model combines three ideas:
\begin{enumerate}
  \item perceived controversy is jointly determined by a propositional and a pragmatic component;
  \item speaker choice depends on a utility trade-off between controversy and persuasive goal strength;
  \item stronger markers incur additional pragmatic costs, which creates a
        \emph{weakest-sufficient} bias and can block overly strong markers.
\end{enumerate}

The three investigated markers are ordered from weakest to strongest as
follows:
\[
\textit{soviel ich wei\ss}
\;<\;
\textit{ja}
\;<\;
\textit{bekanntlich}.
\]

\subsection{Controversy decomposition}

Let
$pc_{\text{prop}} \in [0,1]$ denote propositional controversy and
$pc_{\text{prag}} \in [0,1]$ pragmatic controversy. Overall perceived
controversy is defined as:
\begin{equation}
pc \,=\, PC(pc_{\text{prop}}, pc_{\text{prag}})
\,=\,
pc_{\text{prop}} \cdot pc_{\text{prag}}.
\end{equation}

This multiplicative form captures the idea that a proposition can be globally
contentious while being relatively uncontroversial in the speaker's local
community, and vice versa.

\subsection{Base utility}

Let $g \in [0,1]$ denote speaker goal strength. The speaker's base context
utility is:
\begin{equation}
U(pc,g)
\,=\,
- w_{pc}\,pc
+ w_g\,g
- w_{\mathrm{int}}\,pc\,g,
\label{eq:model-utility}
\end{equation}
where $w_{pc}, w_g, w_{\mathrm{int}} > 0$.

The first term penalizes controversy, the second rewards persuasive effort,
and the interaction term captures that pursuing a strong goal is increasingly
costly when controversy is high.

\subsection{Noisy threshold intervals}

Each marker $u_i$ is associated with a lower threshold $\theta_i$ and an upper
threshold $\theta_{i+1}$, where the upper threshold is the lower threshold of
the next stronger marker. For the strongest marker, the upper boundary is
$+\infty$.

Thresholds are treated as noisy rather than fixed:
\begin{equation}
\theta_i \sim \mathcal{N}(\mu_i, \sigma_i^2),
\end{equation}
with ordered means
\[
\mu_{\textit{soviel ich wei\ss}}
\;<\;
\mu_{\textit{ja}}
\;<\;
\mu_{\textit{bekanntlich}}.
\]

A marker fits the context to the extent that the utility falls inside its
licensed interval. The corresponding noisy interval-fit term is:
\begin{equation}
\mathrm{Fit}_{\mathrm{interval}}(u_i \mid U)
\;\propto\;
\mathbb{E}_{\theta_i}\left[\sigma\!\left(\lambda(U-\theta_i)\right)\right]
\cdot
\mathbb{E}_{\theta_{i+1}}\left[\sigma\!\left(-\lambda(U-\theta_{i+1})\right)\right],
\label{eq:model-fit}
\end{equation}
where $\sigma(x)=(1+e^{-x})^{-1}$ is the logistic function and $\lambda > 0$
controls boundary sharpness.

The first factor rewards utilities above the lower threshold; the second
factor penalizes utilities above the next stronger marker's lower threshold.
Thus, each marker is associated with an interval rather than a single cutoff.

\subsection{Marker costs and sufficiency blocking}

In addition to interval fit, each marker carries a marker-specific pragmatic
cost $c_u \geq 0$. The weakest marker can be treated as the reference level,
with stronger markers incurring larger costs.

Speaker choice is therefore modeled as:
\begin{equation}
P(u \mid pc,g)
\;\propto\;
\mathrm{Fit}_{\mathrm{interval}}\!\left(u \mid U(pc,g)\right)
\cdot e^{-c_u}.
\label{eq:model-choice}
\end{equation}

This cost term introduces a \emph{weakest-sufficient} bias. A stronger marker
may be licensed by the context, but still lose to a weaker marker if the weaker
option already serves the speaker's purpose sufficiently well. This yields the
prediction of \textbf{sufficiency blocking}: in low-controversy contexts, even
under relatively high goal strength, \textit{bekanntlich} can be blocked if
\textit{ja} or even \textit{soviel ich wei\ss} already provides enough force
at lower cost.

\subsection{Listener inversion}

Given an observed marker $u$, the listener inverts the speaker model to infer
latent context variables:
\begin{equation}
P(pc_{\text{prop}}, pc_{\text{prag}}, g \mid u)
\;\propto\;
P\!\left(u \mid PC(pc_{\text{prop}}, pc_{\text{prag}}), g\right)
\,P(pc_{\text{prop}})\,P(pc_{\text{prag}})\,P(g).
\label{eq:model-posterior}
\end{equation}

Because stronger markers generally require higher utility, they shift posterior
mass toward lower controversy and higher goal strength. However, the cost term
weakens a purely monotonic mapping from marker strength to speaker goal:
observing a weaker marker can reflect either lower utility or the fact that a
stronger marker would have been pragmatically excessive.

\subsection{Adoption model}

The listener's downstream adoption decision depends on the posterior over goal
strength and controversy:
\begin{equation}
\text{logit}\left(P_{\mathrm{adopt}}\right)
\;=
\eta_0
+ \eta_g\,\mathbb{E}[g \mid u,pc]
- \eta_{pc}\,pc
- \eta_{\mathrm{int}}\,\mathbb{E}[g \mid u,pc]\,pc.
\label{eq:model-adopt}
\end{equation}

Here, $\eta_0$ is a baseline adoption tendency, $\eta_g$ captures the positive
effect of inferred speaker goal, $\eta_{pc}$ penalizes controversy, and
$\eta_{\mathrm{int}}$ allows controversy to attenuate the effect of inferred
goal strength.

\subsection{Qualitative predictions}

The model yields the following core predictions:
\begin{enumerate}
  \item \textbf{Goal effect:} holding controversy fixed, higher $g$ increases
        the probability of stronger markers.
  \item \textbf{Controversy effect:} holding $g$ fixed, higher $pc$ increases
        the probability of weaker markers.
  \item \textbf{Attenuated goal effect under controversy:} the impact of $g$
        on marker strength is reduced when $pc$ is high.
  \item \textbf{Sufficiency blocking:} in low-$pc$ regions, a stronger marker
        can be blocked by a weaker marker if the weaker one already does enough.
  \item \textbf{Listener-side inversion:} stronger markers tend to induce lower
        inferred controversy and higher inferred goal strength.
  \item \textbf{Credibility discounting:} even when a strong marker supports
        high inferred goal strength, its effect on action adoption is reduced
        in highly controversial contexts.
\end{enumerate}

\subsection{Interpretive summary}

Relative to a vanilla RSA model that directly softmaxes over utterance utility,
this model treats markers as corresponding to \emph{intervals} on a latent
utility scale. The interval-fit component answers the question:
\enquote{How well does this marker match the current utility region?}
The cost term answers the question:
\enquote{Is this marker stronger or more marked than the speaker needs?}
Only their combination determines actual choice.

\end{document}
.
%
% Compile stand-alone:
%   pdflatex model.tex
%
% Embed in main.tex:
%   % model.tex
% Stand-alone description of the cost-augmented noisy-threshold RSA model.
% Can be embedded in main.tex via \input{model}.
%
% Compile stand-alone:
%   pdflatex model.tex
%
% Embed in main.tex:
%   \input{model}

\documentclass[12pt]{article}
\usepackage{CSPArticleStyle}
\usepackage{csquotes}
\usepackage[margin=1in]{geometry}
\usepackage{booktabs}
\usepackage{parskip}
\usepackage{amsmath}
\usepackage{amssymb}
\usepackage{times}
\usepackage[T1]{fontenc}
\usepackage[utf8]{inputenc}
\usepackage{xcolor}

\title{Cost-Augmented Noisy-Threshold RSA Model for Consensus Markers}
\author{}
\date{}

\begin{document}

%% ============================================================
%% When embedding in main.tex, delete from \documentclass to
%% \begin{document} and delete \end{document} at the bottom.
%% ============================================================

\section{Model}

\subsection{Overview}

We model German consensus markers as lexical items that are licensed over
intervals of a latent utility scale. The model combines three ideas:
\begin{enumerate}
  \item perceived controversy is jointly determined by a propositional and a pragmatic component;
  \item speaker choice depends on a utility trade-off between controversy and persuasive goal strength;
  \item stronger markers incur additional pragmatic costs, which creates a
        \emph{weakest-sufficient} bias and can block overly strong markers.
\end{enumerate}

The three investigated markers are ordered from weakest to strongest as
follows:
\[
\textit{soviel ich wei\ss}
\;<\;
\textit{ja}
\;<\;
\textit{bekanntlich}.
\]

\subsection{Controversy decomposition}

Let
$pc_{\text{prop}} \in [0,1]$ denote propositional controversy and
$pc_{\text{prag}} \in [0,1]$ pragmatic controversy. Overall perceived
controversy is defined as:
\begin{equation}
pc \,=\, PC(pc_{\text{prop}}, pc_{\text{prag}})
\,=\,
pc_{\text{prop}} \cdot pc_{\text{prag}}.
\end{equation}

This multiplicative form captures the idea that a proposition can be globally
contentious while being relatively uncontroversial in the speaker's local
community, and vice versa.

\subsection{Base utility}

Let $g \in [0,1]$ denote speaker goal strength. The speaker's base context
utility is:
\begin{equation}
U(pc,g)
\,=\,
- w_{pc}\,pc
+ w_g\,g
- w_{\mathrm{int}}\,pc\,g,
\label{eq:model-utility}
\end{equation}
where $w_{pc}, w_g, w_{\mathrm{int}} > 0$.

The first term penalizes controversy, the second rewards persuasive effort,
and the interaction term captures that pursuing a strong goal is increasingly
costly when controversy is high.

\subsection{Noisy threshold intervals}

Each marker $u_i$ is associated with a lower threshold $\theta_i$ and an upper
threshold $\theta_{i+1}$, where the upper threshold is the lower threshold of
the next stronger marker. For the strongest marker, the upper boundary is
$+\infty$.

Thresholds are treated as noisy rather than fixed:
\begin{equation}
\theta_i \sim \mathcal{N}(\mu_i, \sigma_i^2),
\end{equation}
with ordered means
\[
\mu_{\textit{soviel ich wei\ss}}
\;<\;
\mu_{\textit{ja}}
\;<\;
\mu_{\textit{bekanntlich}}.
\]

A marker fits the context to the extent that the utility falls inside its
licensed interval. The corresponding noisy interval-fit term is:
\begin{equation}
\mathrm{Fit}_{\mathrm{interval}}(u_i \mid U)
\;\propto\;
\mathbb{E}_{\theta_i}\left[\sigma\!\left(\lambda(U-\theta_i)\right)\right]
\cdot
\mathbb{E}_{\theta_{i+1}}\left[\sigma\!\left(-\lambda(U-\theta_{i+1})\right)\right],
\label{eq:model-fit}
\end{equation}
where $\sigma(x)=(1+e^{-x})^{-1}$ is the logistic function and $\lambda > 0$
controls boundary sharpness.

The first factor rewards utilities above the lower threshold; the second
factor penalizes utilities above the next stronger marker's lower threshold.
Thus, each marker is associated with an interval rather than a single cutoff.

\subsection{Marker costs and sufficiency blocking}

In addition to interval fit, each marker carries a marker-specific pragmatic
cost $c_u \geq 0$. The weakest marker can be treated as the reference level,
with stronger markers incurring larger costs.

Speaker choice is therefore modeled as:
\begin{equation}
P(u \mid pc,g)
\;\propto\;
\mathrm{Fit}_{\mathrm{interval}}\!\left(u \mid U(pc,g)\right)
\cdot e^{-c_u}.
\label{eq:model-choice}
\end{equation}

This cost term introduces a \emph{weakest-sufficient} bias. A stronger marker
may be licensed by the context, but still lose to a weaker marker if the weaker
option already serves the speaker's purpose sufficiently well. This yields the
prediction of \textbf{sufficiency blocking}: in low-controversy contexts, even
under relatively high goal strength, \textit{bekanntlich} can be blocked if
\textit{ja} or even \textit{soviel ich wei\ss} already provides enough force
at lower cost.

\subsection{Listener inversion}

Given an observed marker $u$, the listener inverts the speaker model to infer
latent context variables:
\begin{equation}
P(pc_{\text{prop}}, pc_{\text{prag}}, g \mid u)
\;\propto\;
P\!\left(u \mid PC(pc_{\text{prop}}, pc_{\text{prag}}), g\right)
\,P(pc_{\text{prop}})\,P(pc_{\text{prag}})\,P(g).
\label{eq:model-posterior}
\end{equation}

Because stronger markers generally require higher utility, they shift posterior
mass toward lower controversy and higher goal strength. However, the cost term
weakens a purely monotonic mapping from marker strength to speaker goal:
observing a weaker marker can reflect either lower utility or the fact that a
stronger marker would have been pragmatically excessive.

\subsection{Adoption model}

The listener's downstream adoption decision depends on the posterior over goal
strength and controversy:
\begin{equation}
\text{logit}\left(P_{\mathrm{adopt}}\right)
\;=
\eta_0
+ \eta_g\,\mathbb{E}[g \mid u,pc]
- \eta_{pc}\,pc
- \eta_{\mathrm{int}}\,\mathbb{E}[g \mid u,pc]\,pc.
\label{eq:model-adopt}
\end{equation}

Here, $\eta_0$ is a baseline adoption tendency, $\eta_g$ captures the positive
effect of inferred speaker goal, $\eta_{pc}$ penalizes controversy, and
$\eta_{\mathrm{int}}$ allows controversy to attenuate the effect of inferred
goal strength.

\subsection{Qualitative predictions}

The model yields the following core predictions:
\begin{enumerate}
  \item \textbf{Goal effect:} holding controversy fixed, higher $g$ increases
        the probability of stronger markers.
  \item \textbf{Controversy effect:} holding $g$ fixed, higher $pc$ increases
        the probability of weaker markers.
  \item \textbf{Attenuated goal effect under controversy:} the impact of $g$
        on marker strength is reduced when $pc$ is high.
  \item \textbf{Sufficiency blocking:} in low-$pc$ regions, a stronger marker
        can be blocked by a weaker marker if the weaker one already does enough.
  \item \textbf{Listener-side inversion:} stronger markers tend to induce lower
        inferred controversy and higher inferred goal strength.
  \item \textbf{Credibility discounting:} even when a strong marker supports
        high inferred goal strength, its effect on action adoption is reduced
        in highly controversial contexts.
\end{enumerate}

\subsection{Interpretive summary}

Relative to a vanilla RSA model that directly softmaxes over utterance utility,
this model treats markers as corresponding to \emph{intervals} on a latent
utility scale. The interval-fit component answers the question:
\enquote{How well does this marker match the current utility region?}
The cost term answers the question:
\enquote{Is this marker stronger or more marked than the speaker needs?}
Only their combination determines actual choice.

\end{document}


\documentclass[12pt]{article}
\usepackage{CSPArticleStyle}
\usepackage{csquotes}
\usepackage[margin=1in]{geometry}
\usepackage{booktabs}
\usepackage{parskip}
\usepackage{amsmath}
\usepackage{amssymb}
\usepackage{times}
\usepackage[T1]{fontenc}
\usepackage[utf8]{inputenc}
\usepackage{xcolor}

\title{Cost-Augmented Noisy-Threshold RSA Model for Consensus Markers}
\author{}
\date{}

\begin{document}

%% ============================================================
%% When embedding in main.tex, delete from \documentclass to
%% \begin{document} and delete \end{document} at the bottom.
%% ============================================================

\section{Model}

\subsection{Overview}

We model German consensus markers as lexical items that are licensed over
intervals of a latent utility scale. The model combines three ideas:
\begin{enumerate}
  \item perceived controversy is jointly determined by a propositional and a pragmatic component;
  \item speaker choice depends on a utility trade-off between controversy and persuasive goal strength;
  \item stronger markers incur additional pragmatic costs, which creates a
        \emph{weakest-sufficient} bias and can block overly strong markers.
\end{enumerate}

The three investigated markers are ordered from weakest to strongest as
follows:
\[
\textit{soviel ich wei\ss}
\;<\;
\textit{ja}
\;<\;
\textit{bekanntlich}.
\]

\subsection{Controversy decomposition}

Let
$pc_{\text{prop}} \in [0,1]$ denote propositional controversy and
$pc_{\text{prag}} \in [0,1]$ pragmatic controversy. Overall perceived
controversy is defined as:
\begin{equation}
pc \,=\, PC(pc_{\text{prop}}, pc_{\text{prag}})
\,=\,
pc_{\text{prop}} \cdot pc_{\text{prag}}.
\end{equation}

This multiplicative form captures the idea that a proposition can be globally
contentious while being relatively uncontroversial in the speaker's local
community, and vice versa.

\subsection{Base utility}

Let $g \in [0,1]$ denote speaker goal strength. The speaker's base context
utility is:
\begin{equation}
U(pc,g)
\,=\,
- w_{pc}\,pc
+ w_g\,g
- w_{\mathrm{int}}\,pc\,g,
\label{eq:model-utility}
\end{equation}
where $w_{pc}, w_g, w_{\mathrm{int}} > 0$.

The first term penalizes controversy, the second rewards persuasive effort,
and the interaction term captures that pursuing a strong goal is increasingly
costly when controversy is high.

\subsection{Noisy threshold intervals}

Each marker $u_i$ is associated with a lower threshold $\theta_i$ and an upper
threshold $\theta_{i+1}$, where the upper threshold is the lower threshold of
the next stronger marker. For the strongest marker, the upper boundary is
$+\infty$.

Thresholds are treated as noisy rather than fixed:
\begin{equation}
\theta_i \sim \mathcal{N}(\mu_i, \sigma_i^2),
\end{equation}
with ordered means
\[
\mu_{\textit{soviel ich wei\ss}}
\;<\;
\mu_{\textit{ja}}
\;<\;
\mu_{\textit{bekanntlich}}.
\]

A marker fits the context to the extent that the utility falls inside its
licensed interval. The corresponding noisy interval-fit term is:
\begin{equation}
\mathrm{Fit}_{\mathrm{interval}}(u_i \mid U)
\;\propto\;
\mathbb{E}_{\theta_i}\left[\sigma\!\left(\lambda(U-\theta_i)\right)\right]
\cdot
\mathbb{E}_{\theta_{i+1}}\left[\sigma\!\left(-\lambda(U-\theta_{i+1})\right)\right],
\label{eq:model-fit}
\end{equation}
where $\sigma(x)=(1+e^{-x})^{-1}$ is the logistic function and $\lambda > 0$
controls boundary sharpness.

The first factor rewards utilities above the lower threshold; the second
factor penalizes utilities above the next stronger marker's lower threshold.
Thus, each marker is associated with an interval rather than a single cutoff.

\subsection{Marker costs and sufficiency blocking}

In addition to interval fit, each marker carries a marker-specific pragmatic
cost $c_u \geq 0$. The weakest marker can be treated as the reference level,
with stronger markers incurring larger costs.

Speaker choice is therefore modeled as:
\begin{equation}
P(u \mid pc,g)
\;\propto\;
\mathrm{Fit}_{\mathrm{interval}}\!\left(u \mid U(pc,g)\right)
\cdot e^{-c_u}.
\label{eq:model-choice}
\end{equation}

This cost term introduces a \emph{weakest-sufficient} bias. A stronger marker
may be licensed by the context, but still lose to a weaker marker if the weaker
option already serves the speaker's purpose sufficiently well. This yields the
prediction of \textbf{sufficiency blocking}: in low-controversy contexts, even
under relatively high goal strength, \textit{bekanntlich} can be blocked if
\textit{ja} or even \textit{soviel ich wei\ss} already provides enough force
at lower cost.

\subsection{Listener inversion}

Given an observed marker $u$, the listener inverts the speaker model to infer
latent context variables:
\begin{equation}
P(pc_{\text{prop}}, pc_{\text{prag}}, g \mid u)
\;\propto\;
P\!\left(u \mid PC(pc_{\text{prop}}, pc_{\text{prag}}), g\right)
\,P(pc_{\text{prop}})\,P(pc_{\text{prag}})\,P(g).
\label{eq:model-posterior}
\end{equation}

Because stronger markers generally require higher utility, they shift posterior
mass toward lower controversy and higher goal strength. However, the cost term
weakens a purely monotonic mapping from marker strength to speaker goal:
observing a weaker marker can reflect either lower utility or the fact that a
stronger marker would have been pragmatically excessive.

\subsection{Adoption model}

The listener's downstream adoption decision depends on the posterior over goal
strength and controversy:
\begin{equation}
\text{logit}\left(P_{\mathrm{adopt}}\right)
\;=
\eta_0
+ \eta_g\,\mathbb{E}[g \mid u,pc]
- \eta_{pc}\,pc
- \eta_{\mathrm{int}}\,\mathbb{E}[g \mid u,pc]\,pc.
\label{eq:model-adopt}
\end{equation}

Here, $\eta_0$ is a baseline adoption tendency, $\eta_g$ captures the positive
effect of inferred speaker goal, $\eta_{pc}$ penalizes controversy, and
$\eta_{\mathrm{int}}$ allows controversy to attenuate the effect of inferred
goal strength.

\subsection{Qualitative predictions}

The model yields the following core predictions:
\begin{enumerate}
  \item \textbf{Goal effect:} holding controversy fixed, higher $g$ increases
        the probability of stronger markers.
  \item \textbf{Controversy effect:} holding $g$ fixed, higher $pc$ increases
        the probability of weaker markers.
  \item \textbf{Attenuated goal effect under controversy:} the impact of $g$
        on marker strength is reduced when $pc$ is high.
  \item \textbf{Sufficiency blocking:} in low-$pc$ regions, a stronger marker
        can be blocked by a weaker marker if the weaker one already does enough.
  \item \textbf{Listener-side inversion:} stronger markers tend to induce lower
        inferred controversy and higher inferred goal strength.
  \item \textbf{Credibility discounting:} even when a strong marker supports
        high inferred goal strength, its effect on action adoption is reduced
        in highly controversial contexts.
\end{enumerate}

\subsection{Interpretive summary}

Relative to a vanilla RSA model that directly softmaxes over utterance utility,
this model treats markers as corresponding to \emph{intervals} on a latent
utility scale. The interval-fit component answers the question:
\enquote{How well does this marker match the current utility region?}
The cost term answers the question:
\enquote{Is this marker stronger or more marked than the speaker needs?}
Only their combination determines actual choice.

\end{document}
.
%
% Compile stand-alone:
%   pdflatex model.tex
%
% Embed in main.tex:
%   % model.tex
% Stand-alone description of the cost-augmented noisy-threshold RSA model.
% Can be embedded in main.tex via % model.tex
% Stand-alone description of the cost-augmented noisy-threshold RSA model.
% Can be embedded in main.tex via \input{model}.
%
% Compile stand-alone:
%   pdflatex model.tex
%
% Embed in main.tex:
%   \input{model}

\documentclass[12pt]{article}
\usepackage{CSPArticleStyle}
\usepackage{csquotes}
\usepackage[margin=1in]{geometry}
\usepackage{booktabs}
\usepackage{parskip}
\usepackage{amsmath}
\usepackage{amssymb}
\usepackage{times}
\usepackage[T1]{fontenc}
\usepackage[utf8]{inputenc}
\usepackage{xcolor}

\title{Cost-Augmented Noisy-Threshold RSA Model for Consensus Markers}
\author{}
\date{}

\begin{document}

%% ============================================================
%% When embedding in main.tex, delete from \documentclass to
%% \begin{document} and delete \end{document} at the bottom.
%% ============================================================

\section{Model}

\subsection{Overview}

We model German consensus markers as lexical items that are licensed over
intervals of a latent utility scale. The model combines three ideas:
\begin{enumerate}
  \item perceived controversy is jointly determined by a propositional and a pragmatic component;
  \item speaker choice depends on a utility trade-off between controversy and persuasive goal strength;
  \item stronger markers incur additional pragmatic costs, which creates a
        \emph{weakest-sufficient} bias and can block overly strong markers.
\end{enumerate}

The three investigated markers are ordered from weakest to strongest as
follows:
\[
\textit{soviel ich wei\ss}
\;<\;
\textit{ja}
\;<\;
\textit{bekanntlich}.
\]

\subsection{Controversy decomposition}

Let
$pc_{\text{prop}} \in [0,1]$ denote propositional controversy and
$pc_{\text{prag}} \in [0,1]$ pragmatic controversy. Overall perceived
controversy is defined as:
\begin{equation}
pc \,=\, PC(pc_{\text{prop}}, pc_{\text{prag}})
\,=\,
pc_{\text{prop}} \cdot pc_{\text{prag}}.
\end{equation}

This multiplicative form captures the idea that a proposition can be globally
contentious while being relatively uncontroversial in the speaker's local
community, and vice versa.

\subsection{Base utility}

Let $g \in [0,1]$ denote speaker goal strength. The speaker's base context
utility is:
\begin{equation}
U(pc,g)
\,=\,
- w_{pc}\,pc
+ w_g\,g
- w_{\mathrm{int}}\,pc\,g,
\label{eq:model-utility}
\end{equation}
where $w_{pc}, w_g, w_{\mathrm{int}} > 0$.

The first term penalizes controversy, the second rewards persuasive effort,
and the interaction term captures that pursuing a strong goal is increasingly
costly when controversy is high.

\subsection{Noisy threshold intervals}

Each marker $u_i$ is associated with a lower threshold $\theta_i$ and an upper
threshold $\theta_{i+1}$, where the upper threshold is the lower threshold of
the next stronger marker. For the strongest marker, the upper boundary is
$+\infty$.

Thresholds are treated as noisy rather than fixed:
\begin{equation}
\theta_i \sim \mathcal{N}(\mu_i, \sigma_i^2),
\end{equation}
with ordered means
\[
\mu_{\textit{soviel ich wei\ss}}
\;<\;
\mu_{\textit{ja}}
\;<\;
\mu_{\textit{bekanntlich}}.
\]

A marker fits the context to the extent that the utility falls inside its
licensed interval. The corresponding noisy interval-fit term is:
\begin{equation}
\mathrm{Fit}_{\mathrm{interval}}(u_i \mid U)
\;\propto\;
\mathbb{E}_{\theta_i}\left[\sigma\!\left(\lambda(U-\theta_i)\right)\right]
\cdot
\mathbb{E}_{\theta_{i+1}}\left[\sigma\!\left(-\lambda(U-\theta_{i+1})\right)\right],
\label{eq:model-fit}
\end{equation}
where $\sigma(x)=(1+e^{-x})^{-1}$ is the logistic function and $\lambda > 0$
controls boundary sharpness.

The first factor rewards utilities above the lower threshold; the second
factor penalizes utilities above the next stronger marker's lower threshold.
Thus, each marker is associated with an interval rather than a single cutoff.

\subsection{Marker costs and sufficiency blocking}

In addition to interval fit, each marker carries a marker-specific pragmatic
cost $c_u \geq 0$. The weakest marker can be treated as the reference level,
with stronger markers incurring larger costs.

Speaker choice is therefore modeled as:
\begin{equation}
P(u \mid pc,g)
\;\propto\;
\mathrm{Fit}_{\mathrm{interval}}\!\left(u \mid U(pc,g)\right)
\cdot e^{-c_u}.
\label{eq:model-choice}
\end{equation}

This cost term introduces a \emph{weakest-sufficient} bias. A stronger marker
may be licensed by the context, but still lose to a weaker marker if the weaker
option already serves the speaker's purpose sufficiently well. This yields the
prediction of \textbf{sufficiency blocking}: in low-controversy contexts, even
under relatively high goal strength, \textit{bekanntlich} can be blocked if
\textit{ja} or even \textit{soviel ich wei\ss} already provides enough force
at lower cost.

\subsection{Listener inversion}

Given an observed marker $u$, the listener inverts the speaker model to infer
latent context variables:
\begin{equation}
P(pc_{\text{prop}}, pc_{\text{prag}}, g \mid u)
\;\propto\;
P\!\left(u \mid PC(pc_{\text{prop}}, pc_{\text{prag}}), g\right)
\,P(pc_{\text{prop}})\,P(pc_{\text{prag}})\,P(g).
\label{eq:model-posterior}
\end{equation}

Because stronger markers generally require higher utility, they shift posterior
mass toward lower controversy and higher goal strength. However, the cost term
weakens a purely monotonic mapping from marker strength to speaker goal:
observing a weaker marker can reflect either lower utility or the fact that a
stronger marker would have been pragmatically excessive.

\subsection{Adoption model}

The listener's downstream adoption decision depends on the posterior over goal
strength and controversy:
\begin{equation}
\text{logit}\left(P_{\mathrm{adopt}}\right)
\;=
\eta_0
+ \eta_g\,\mathbb{E}[g \mid u,pc]
- \eta_{pc}\,pc
- \eta_{\mathrm{int}}\,\mathbb{E}[g \mid u,pc]\,pc.
\label{eq:model-adopt}
\end{equation}

Here, $\eta_0$ is a baseline adoption tendency, $\eta_g$ captures the positive
effect of inferred speaker goal, $\eta_{pc}$ penalizes controversy, and
$\eta_{\mathrm{int}}$ allows controversy to attenuate the effect of inferred
goal strength.

\subsection{Qualitative predictions}

The model yields the following core predictions:
\begin{enumerate}
  \item \textbf{Goal effect:} holding controversy fixed, higher $g$ increases
        the probability of stronger markers.
  \item \textbf{Controversy effect:} holding $g$ fixed, higher $pc$ increases
        the probability of weaker markers.
  \item \textbf{Attenuated goal effect under controversy:} the impact of $g$
        on marker strength is reduced when $pc$ is high.
  \item \textbf{Sufficiency blocking:} in low-$pc$ regions, a stronger marker
        can be blocked by a weaker marker if the weaker one already does enough.
  \item \textbf{Listener-side inversion:} stronger markers tend to induce lower
        inferred controversy and higher inferred goal strength.
  \item \textbf{Credibility discounting:} even when a strong marker supports
        high inferred goal strength, its effect on action adoption is reduced
        in highly controversial contexts.
\end{enumerate}

\subsection{Interpretive summary}

Relative to a vanilla RSA model that directly softmaxes over utterance utility,
this model treats markers as corresponding to \emph{intervals} on a latent
utility scale. The interval-fit component answers the question:
\enquote{How well does this marker match the current utility region?}
The cost term answers the question:
\enquote{Is this marker stronger or more marked than the speaker needs?}
Only their combination determines actual choice.

\end{document}
.
%
% Compile stand-alone:
%   pdflatex model.tex
%
% Embed in main.tex:
%   % model.tex
% Stand-alone description of the cost-augmented noisy-threshold RSA model.
% Can be embedded in main.tex via \input{model}.
%
% Compile stand-alone:
%   pdflatex model.tex
%
% Embed in main.tex:
%   \input{model}

\documentclass[12pt]{article}
\usepackage{CSPArticleStyle}
\usepackage{csquotes}
\usepackage[margin=1in]{geometry}
\usepackage{booktabs}
\usepackage{parskip}
\usepackage{amsmath}
\usepackage{amssymb}
\usepackage{times}
\usepackage[T1]{fontenc}
\usepackage[utf8]{inputenc}
\usepackage{xcolor}

\title{Cost-Augmented Noisy-Threshold RSA Model for Consensus Markers}
\author{}
\date{}

\begin{document}

%% ============================================================
%% When embedding in main.tex, delete from \documentclass to
%% \begin{document} and delete \end{document} at the bottom.
%% ============================================================

\section{Model}

\subsection{Overview}

We model German consensus markers as lexical items that are licensed over
intervals of a latent utility scale. The model combines three ideas:
\begin{enumerate}
  \item perceived controversy is jointly determined by a propositional and a pragmatic component;
  \item speaker choice depends on a utility trade-off between controversy and persuasive goal strength;
  \item stronger markers incur additional pragmatic costs, which creates a
        \emph{weakest-sufficient} bias and can block overly strong markers.
\end{enumerate}

The three investigated markers are ordered from weakest to strongest as
follows:
\[
\textit{soviel ich wei\ss}
\;<\;
\textit{ja}
\;<\;
\textit{bekanntlich}.
\]

\subsection{Controversy decomposition}

Let
$pc_{\text{prop}} \in [0,1]$ denote propositional controversy and
$pc_{\text{prag}} \in [0,1]$ pragmatic controversy. Overall perceived
controversy is defined as:
\begin{equation}
pc \,=\, PC(pc_{\text{prop}}, pc_{\text{prag}})
\,=\,
pc_{\text{prop}} \cdot pc_{\text{prag}}.
\end{equation}

This multiplicative form captures the idea that a proposition can be globally
contentious while being relatively uncontroversial in the speaker's local
community, and vice versa.

\subsection{Base utility}

Let $g \in [0,1]$ denote speaker goal strength. The speaker's base context
utility is:
\begin{equation}
U(pc,g)
\,=\,
- w_{pc}\,pc
+ w_g\,g
- w_{\mathrm{int}}\,pc\,g,
\label{eq:model-utility}
\end{equation}
where $w_{pc}, w_g, w_{\mathrm{int}} > 0$.

The first term penalizes controversy, the second rewards persuasive effort,
and the interaction term captures that pursuing a strong goal is increasingly
costly when controversy is high.

\subsection{Noisy threshold intervals}

Each marker $u_i$ is associated with a lower threshold $\theta_i$ and an upper
threshold $\theta_{i+1}$, where the upper threshold is the lower threshold of
the next stronger marker. For the strongest marker, the upper boundary is
$+\infty$.

Thresholds are treated as noisy rather than fixed:
\begin{equation}
\theta_i \sim \mathcal{N}(\mu_i, \sigma_i^2),
\end{equation}
with ordered means
\[
\mu_{\textit{soviel ich wei\ss}}
\;<\;
\mu_{\textit{ja}}
\;<\;
\mu_{\textit{bekanntlich}}.
\]

A marker fits the context to the extent that the utility falls inside its
licensed interval. The corresponding noisy interval-fit term is:
\begin{equation}
\mathrm{Fit}_{\mathrm{interval}}(u_i \mid U)
\;\propto\;
\mathbb{E}_{\theta_i}\left[\sigma\!\left(\lambda(U-\theta_i)\right)\right]
\cdot
\mathbb{E}_{\theta_{i+1}}\left[\sigma\!\left(-\lambda(U-\theta_{i+1})\right)\right],
\label{eq:model-fit}
\end{equation}
where $\sigma(x)=(1+e^{-x})^{-1}$ is the logistic function and $\lambda > 0$
controls boundary sharpness.

The first factor rewards utilities above the lower threshold; the second
factor penalizes utilities above the next stronger marker's lower threshold.
Thus, each marker is associated with an interval rather than a single cutoff.

\subsection{Marker costs and sufficiency blocking}

In addition to interval fit, each marker carries a marker-specific pragmatic
cost $c_u \geq 0$. The weakest marker can be treated as the reference level,
with stronger markers incurring larger costs.

Speaker choice is therefore modeled as:
\begin{equation}
P(u \mid pc,g)
\;\propto\;
\mathrm{Fit}_{\mathrm{interval}}\!\left(u \mid U(pc,g)\right)
\cdot e^{-c_u}.
\label{eq:model-choice}
\end{equation}

This cost term introduces a \emph{weakest-sufficient} bias. A stronger marker
may be licensed by the context, but still lose to a weaker marker if the weaker
option already serves the speaker's purpose sufficiently well. This yields the
prediction of \textbf{sufficiency blocking}: in low-controversy contexts, even
under relatively high goal strength, \textit{bekanntlich} can be blocked if
\textit{ja} or even \textit{soviel ich wei\ss} already provides enough force
at lower cost.

\subsection{Listener inversion}

Given an observed marker $u$, the listener inverts the speaker model to infer
latent context variables:
\begin{equation}
P(pc_{\text{prop}}, pc_{\text{prag}}, g \mid u)
\;\propto\;
P\!\left(u \mid PC(pc_{\text{prop}}, pc_{\text{prag}}), g\right)
\,P(pc_{\text{prop}})\,P(pc_{\text{prag}})\,P(g).
\label{eq:model-posterior}
\end{equation}

Because stronger markers generally require higher utility, they shift posterior
mass toward lower controversy and higher goal strength. However, the cost term
weakens a purely monotonic mapping from marker strength to speaker goal:
observing a weaker marker can reflect either lower utility or the fact that a
stronger marker would have been pragmatically excessive.

\subsection{Adoption model}

The listener's downstream adoption decision depends on the posterior over goal
strength and controversy:
\begin{equation}
\text{logit}\left(P_{\mathrm{adopt}}\right)
\;=
\eta_0
+ \eta_g\,\mathbb{E}[g \mid u,pc]
- \eta_{pc}\,pc
- \eta_{\mathrm{int}}\,\mathbb{E}[g \mid u,pc]\,pc.
\label{eq:model-adopt}
\end{equation}

Here, $\eta_0$ is a baseline adoption tendency, $\eta_g$ captures the positive
effect of inferred speaker goal, $\eta_{pc}$ penalizes controversy, and
$\eta_{\mathrm{int}}$ allows controversy to attenuate the effect of inferred
goal strength.

\subsection{Qualitative predictions}

The model yields the following core predictions:
\begin{enumerate}
  \item \textbf{Goal effect:} holding controversy fixed, higher $g$ increases
        the probability of stronger markers.
  \item \textbf{Controversy effect:} holding $g$ fixed, higher $pc$ increases
        the probability of weaker markers.
  \item \textbf{Attenuated goal effect under controversy:} the impact of $g$
        on marker strength is reduced when $pc$ is high.
  \item \textbf{Sufficiency blocking:} in low-$pc$ regions, a stronger marker
        can be blocked by a weaker marker if the weaker one already does enough.
  \item \textbf{Listener-side inversion:} stronger markers tend to induce lower
        inferred controversy and higher inferred goal strength.
  \item \textbf{Credibility discounting:} even when a strong marker supports
        high inferred goal strength, its effect on action adoption is reduced
        in highly controversial contexts.
\end{enumerate}

\subsection{Interpretive summary}

Relative to a vanilla RSA model that directly softmaxes over utterance utility,
this model treats markers as corresponding to \emph{intervals} on a latent
utility scale. The interval-fit component answers the question:
\enquote{How well does this marker match the current utility region?}
The cost term answers the question:
\enquote{Is this marker stronger or more marked than the speaker needs?}
Only their combination determines actual choice.

\end{document}


\documentclass[12pt]{article}
\usepackage{CSPArticleStyle}
\usepackage{csquotes}
\usepackage[margin=1in]{geometry}
\usepackage{booktabs}
\usepackage{parskip}
\usepackage{amsmath}
\usepackage{amssymb}
\usepackage{times}
\usepackage[T1]{fontenc}
\usepackage[utf8]{inputenc}
\usepackage{xcolor}

\title{Cost-Augmented Noisy-Threshold RSA Model for Consensus Markers}
\author{}
\date{}

\begin{document}

%% ============================================================
%% When embedding in main.tex, delete from \documentclass to
%% \begin{document} and delete \end{document} at the bottom.
%% ============================================================

\section{Model}

\subsection{Overview}

We model German consensus markers as lexical items that are licensed over
intervals of a latent utility scale. The model combines three ideas:
\begin{enumerate}
  \item perceived controversy is jointly determined by a propositional and a pragmatic component;
  \item speaker choice depends on a utility trade-off between controversy and persuasive goal strength;
  \item stronger markers incur additional pragmatic costs, which creates a
        \emph{weakest-sufficient} bias and can block overly strong markers.
\end{enumerate}

The three investigated markers are ordered from weakest to strongest as
follows:
\[
\textit{soviel ich wei\ss}
\;<\;
\textit{ja}
\;<\;
\textit{bekanntlich}.
\]

\subsection{Controversy decomposition}

Let
$pc_{\text{prop}} \in [0,1]$ denote propositional controversy and
$pc_{\text{prag}} \in [0,1]$ pragmatic controversy. Overall perceived
controversy is defined as:
\begin{equation}
pc \,=\, PC(pc_{\text{prop}}, pc_{\text{prag}})
\,=\,
pc_{\text{prop}} \cdot pc_{\text{prag}}.
\end{equation}

This multiplicative form captures the idea that a proposition can be globally
contentious while being relatively uncontroversial in the speaker's local
community, and vice versa.

\subsection{Base utility}

Let $g \in [0,1]$ denote speaker goal strength. The speaker's base context
utility is:
\begin{equation}
U(pc,g)
\,=\,
- w_{pc}\,pc
+ w_g\,g
- w_{\mathrm{int}}\,pc\,g,
\label{eq:model-utility}
\end{equation}
where $w_{pc}, w_g, w_{\mathrm{int}} > 0$.

The first term penalizes controversy, the second rewards persuasive effort,
and the interaction term captures that pursuing a strong goal is increasingly
costly when controversy is high.

\subsection{Noisy threshold intervals}

Each marker $u_i$ is associated with a lower threshold $\theta_i$ and an upper
threshold $\theta_{i+1}$, where the upper threshold is the lower threshold of
the next stronger marker. For the strongest marker, the upper boundary is
$+\infty$.

Thresholds are treated as noisy rather than fixed:
\begin{equation}
\theta_i \sim \mathcal{N}(\mu_i, \sigma_i^2),
\end{equation}
with ordered means
\[
\mu_{\textit{soviel ich wei\ss}}
\;<\;
\mu_{\textit{ja}}
\;<\;
\mu_{\textit{bekanntlich}}.
\]

A marker fits the context to the extent that the utility falls inside its
licensed interval. The corresponding noisy interval-fit term is:
\begin{equation}
\mathrm{Fit}_{\mathrm{interval}}(u_i \mid U)
\;\propto\;
\mathbb{E}_{\theta_i}\left[\sigma\!\left(\lambda(U-\theta_i)\right)\right]
\cdot
\mathbb{E}_{\theta_{i+1}}\left[\sigma\!\left(-\lambda(U-\theta_{i+1})\right)\right],
\label{eq:model-fit}
\end{equation}
where $\sigma(x)=(1+e^{-x})^{-1}$ is the logistic function and $\lambda > 0$
controls boundary sharpness.

The first factor rewards utilities above the lower threshold; the second
factor penalizes utilities above the next stronger marker's lower threshold.
Thus, each marker is associated with an interval rather than a single cutoff.

\subsection{Marker costs and sufficiency blocking}

In addition to interval fit, each marker carries a marker-specific pragmatic
cost $c_u \geq 0$. The weakest marker can be treated as the reference level,
with stronger markers incurring larger costs.

Speaker choice is therefore modeled as:
\begin{equation}
P(u \mid pc,g)
\;\propto\;
\mathrm{Fit}_{\mathrm{interval}}\!\left(u \mid U(pc,g)\right)
\cdot e^{-c_u}.
\label{eq:model-choice}
\end{equation}

This cost term introduces a \emph{weakest-sufficient} bias. A stronger marker
may be licensed by the context, but still lose to a weaker marker if the weaker
option already serves the speaker's purpose sufficiently well. This yields the
prediction of \textbf{sufficiency blocking}: in low-controversy contexts, even
under relatively high goal strength, \textit{bekanntlich} can be blocked if
\textit{ja} or even \textit{soviel ich wei\ss} already provides enough force
at lower cost.

\subsection{Listener inversion}

Given an observed marker $u$, the listener inverts the speaker model to infer
latent context variables:
\begin{equation}
P(pc_{\text{prop}}, pc_{\text{prag}}, g \mid u)
\;\propto\;
P\!\left(u \mid PC(pc_{\text{prop}}, pc_{\text{prag}}), g\right)
\,P(pc_{\text{prop}})\,P(pc_{\text{prag}})\,P(g).
\label{eq:model-posterior}
\end{equation}

Because stronger markers generally require higher utility, they shift posterior
mass toward lower controversy and higher goal strength. However, the cost term
weakens a purely monotonic mapping from marker strength to speaker goal:
observing a weaker marker can reflect either lower utility or the fact that a
stronger marker would have been pragmatically excessive.

\subsection{Adoption model}

The listener's downstream adoption decision depends on the posterior over goal
strength and controversy:
\begin{equation}
\text{logit}\left(P_{\mathrm{adopt}}\right)
\;=
\eta_0
+ \eta_g\,\mathbb{E}[g \mid u,pc]
- \eta_{pc}\,pc
- \eta_{\mathrm{int}}\,\mathbb{E}[g \mid u,pc]\,pc.
\label{eq:model-adopt}
\end{equation}

Here, $\eta_0$ is a baseline adoption tendency, $\eta_g$ captures the positive
effect of inferred speaker goal, $\eta_{pc}$ penalizes controversy, and
$\eta_{\mathrm{int}}$ allows controversy to attenuate the effect of inferred
goal strength.

\subsection{Qualitative predictions}

The model yields the following core predictions:
\begin{enumerate}
  \item \textbf{Goal effect:} holding controversy fixed, higher $g$ increases
        the probability of stronger markers.
  \item \textbf{Controversy effect:} holding $g$ fixed, higher $pc$ increases
        the probability of weaker markers.
  \item \textbf{Attenuated goal effect under controversy:} the impact of $g$
        on marker strength is reduced when $pc$ is high.
  \item \textbf{Sufficiency blocking:} in low-$pc$ regions, a stronger marker
        can be blocked by a weaker marker if the weaker one already does enough.
  \item \textbf{Listener-side inversion:} stronger markers tend to induce lower
        inferred controversy and higher inferred goal strength.
  \item \textbf{Credibility discounting:} even when a strong marker supports
        high inferred goal strength, its effect on action adoption is reduced
        in highly controversial contexts.
\end{enumerate}

\subsection{Interpretive summary}

Relative to a vanilla RSA model that directly softmaxes over utterance utility,
this model treats markers as corresponding to \emph{intervals} on a latent
utility scale. The interval-fit component answers the question:
\enquote{How well does this marker match the current utility region?}
The cost term answers the question:
\enquote{Is this marker stronger or more marked than the speaker needs?}
Only their combination determines actual choice.

\end{document}


\documentclass[12pt]{article}
\usepackage{CSPArticleStyle}
\usepackage{csquotes}
\usepackage[margin=1in]{geometry}
\usepackage{booktabs}
\usepackage{parskip}
\usepackage{amsmath}
\usepackage{amssymb}
\usepackage{times}
\usepackage[T1]{fontenc}
\usepackage[utf8]{inputenc}
\usepackage{xcolor}

\title{Cost-Augmented Noisy-Threshold RSA Model for Consensus Markers}
\author{}
\date{}

\begin{document}

%% ============================================================
%% When embedding in main.tex, delete from \documentclass to
%% \begin{document} and delete \end{document} at the bottom.
%% ============================================================

\section{Model}

\subsection{Overview}

We model German consensus markers as lexical items that are licensed over
intervals of a latent utility scale. The model combines three ideas:
\begin{enumerate}
  \item perceived controversy is jointly determined by a propositional and a pragmatic component;
  \item speaker choice depends on a utility trade-off between controversy and persuasive goal strength;
  \item stronger markers incur additional pragmatic costs, which creates a
        \emph{weakest-sufficient} bias and can block overly strong markers.
\end{enumerate}

The three investigated markers are ordered from weakest to strongest as
follows:
\[
\textit{soviel ich wei\ss}
\;<\;
\textit{ja}
\;<\;
\textit{bekanntlich}.
\]

\subsection{Controversy decomposition}

Let
$pc_{\text{prop}} \in [0,1]$ denote propositional controversy and
$pc_{\text{prag}} \in [0,1]$ pragmatic controversy. Overall perceived
controversy is defined as:
\begin{equation}
pc \,=\, PC(pc_{\text{prop}}, pc_{\text{prag}})
\,=\,
pc_{\text{prop}} \cdot pc_{\text{prag}}.
\end{equation}

This multiplicative form captures the idea that a proposition can be globally
contentious while being relatively uncontroversial in the speaker's local
community, and vice versa.

\subsection{Base utility}

Let $g \in [0,1]$ denote speaker goal strength. The speaker's base context
utility is:
\begin{equation}
U(pc,g)
\,=\,
- w_{pc}\,pc
+ w_g\,g
- w_{\mathrm{int}}\,pc\,g,
\label{eq:model-utility}
\end{equation}
where $w_{pc}, w_g, w_{\mathrm{int}} > 0$.

The first term penalizes controversy, the second rewards persuasive effort,
and the interaction term captures that pursuing a strong goal is increasingly
costly when controversy is high.

\subsection{Noisy threshold intervals}

Each marker $u_i$ is associated with a lower threshold $\theta_i$ and an upper
threshold $\theta_{i+1}$, where the upper threshold is the lower threshold of
the next stronger marker. For the strongest marker, the upper boundary is
$+\infty$.

Thresholds are treated as noisy rather than fixed:
\begin{equation}
\theta_i \sim \mathcal{N}(\mu_i, \sigma_i^2),
\end{equation}
with ordered means
\[
\mu_{\textit{soviel ich wei\ss}}
\;<\;
\mu_{\textit{ja}}
\;<\;
\mu_{\textit{bekanntlich}}.
\]

A marker fits the context to the extent that the utility falls inside its
licensed interval. The corresponding noisy interval-fit term is:
\begin{equation}
\mathrm{Fit}_{\mathrm{interval}}(u_i \mid U)
\;\propto\;
\mathbb{E}_{\theta_i}\left[\sigma\!\left(\lambda(U-\theta_i)\right)\right]
\cdot
\mathbb{E}_{\theta_{i+1}}\left[\sigma\!\left(-\lambda(U-\theta_{i+1})\right)\right],
\label{eq:model-fit}
\end{equation}
where $\sigma(x)=(1+e^{-x})^{-1}$ is the logistic function and $\lambda > 0$
controls boundary sharpness.

The first factor rewards utilities above the lower threshold; the second
factor penalizes utilities above the next stronger marker's lower threshold.
Thus, each marker is associated with an interval rather than a single cutoff.

\subsection{Marker costs and sufficiency blocking}

In addition to interval fit, each marker carries a marker-specific pragmatic
cost $c_u \geq 0$. The weakest marker can be treated as the reference level,
with stronger markers incurring larger costs.

Speaker choice is therefore modeled as:
\begin{equation}
P(u \mid pc,g)
\;\propto\;
\mathrm{Fit}_{\mathrm{interval}}\!\left(u \mid U(pc,g)\right)
\cdot e^{-c_u}.
\label{eq:model-choice}
\end{equation}

This cost term introduces a \emph{weakest-sufficient} bias. A stronger marker
may be licensed by the context, but still lose to a weaker marker if the weaker
option already serves the speaker's purpose sufficiently well. This yields the
prediction of \textbf{sufficiency blocking}: in low-controversy contexts, even
under relatively high goal strength, \textit{bekanntlich} can be blocked if
\textit{ja} or even \textit{soviel ich wei\ss} already provides enough force
at lower cost.

\subsection{Listener inversion}

Given an observed marker $u$, the listener inverts the speaker model to infer
latent context variables:
\begin{equation}
P(pc_{\text{prop}}, pc_{\text{prag}}, g \mid u)
\;\propto\;
P\!\left(u \mid PC(pc_{\text{prop}}, pc_{\text{prag}}), g\right)
\,P(pc_{\text{prop}})\,P(pc_{\text{prag}})\,P(g).
\label{eq:model-posterior}
\end{equation}

Because stronger markers generally require higher utility, they shift posterior
mass toward lower controversy and higher goal strength. However, the cost term
weakens a purely monotonic mapping from marker strength to speaker goal:
observing a weaker marker can reflect either lower utility or the fact that a
stronger marker would have been pragmatically excessive.

\subsection{Adoption model}

The listener's downstream adoption decision depends on the posterior over goal
strength and controversy:
\begin{equation}
\text{logit}\left(P_{\mathrm{adopt}}\right)
\;=
\eta_0
+ \eta_g\,\mathbb{E}[g \mid u,pc]
- \eta_{pc}\,pc
- \eta_{\mathrm{int}}\,\mathbb{E}[g \mid u,pc]\,pc.
\label{eq:model-adopt}
\end{equation}

Here, $\eta_0$ is a baseline adoption tendency, $\eta_g$ captures the positive
effect of inferred speaker goal, $\eta_{pc}$ penalizes controversy, and
$\eta_{\mathrm{int}}$ allows controversy to attenuate the effect of inferred
goal strength.

\subsection{Qualitative predictions}

The model yields the following core predictions:
\begin{enumerate}
  \item \textbf{Goal effect:} holding controversy fixed, higher $g$ increases
        the probability of stronger markers.
  \item \textbf{Controversy effect:} holding $g$ fixed, higher $pc$ increases
        the probability of weaker markers.
  \item \textbf{Attenuated goal effect under controversy:} the impact of $g$
        on marker strength is reduced when $pc$ is high.
  \item \textbf{Sufficiency blocking:} in low-$pc$ regions, a stronger marker
        can be blocked by a weaker marker if the weaker one already does enough.
  \item \textbf{Listener-side inversion:} stronger markers tend to induce lower
        inferred controversy and higher inferred goal strength.
  \item \textbf{Credibility discounting:} even when a strong marker supports
        high inferred goal strength, its effect on action adoption is reduced
        in highly controversial contexts.
\end{enumerate}

\subsection{Interpretive summary}

Relative to a vanilla RSA model that directly softmaxes over utterance utility,
this model treats markers as corresponding to \emph{intervals} on a latent
utility scale. The interval-fit component answers the question:
\enquote{How well does this marker match the current utility region?}
The cost term answers the question:
\enquote{Is this marker stronger or more marked than the speaker needs?}
Only their combination determines actual choice.

\end{document}


\documentclass[12pt]{article}
\usepackage{CSPArticleStyle}
\usepackage{csquotes}
\usepackage[margin=1in]{geometry}
\usepackage{booktabs}
\usepackage{parskip}
\usepackage{amsmath}
\usepackage{amssymb}
\usepackage{times}
\usepackage[T1]{fontenc}
\usepackage[utf8]{inputenc}
\usepackage{xcolor}

\title{Cost-Augmented Noisy-Threshold RSA Model for Consensus Markers}
\author{}
\date{}

\begin{document}

%% ============================================================
%% When embedding in main.tex, delete from \documentclass to
%% \begin{document} and delete \end{document} at the bottom.
%% ============================================================

\section{Model}

\subsection{Overview}

We model German consensus markers as lexical items that are licensed over
intervals of a latent utility scale. The model combines three ideas:
\begin{enumerate}
  \item perceived controversy is jointly determined by a propositional and a pragmatic component;
  \item speaker choice depends on a utility trade-off between controversy and persuasive goal strength;
  \item stronger markers incur additional pragmatic costs, which creates a
        \emph{weakest-sufficient} bias and can block overly strong markers.
\end{enumerate}

The three investigated markers are ordered from weakest to strongest as
follows:
\[
\textit{soviel ich wei\ss}
\;<\;
\textit{ja}
\;<\;
\textit{bekanntlich}.
\]

\subsection{Controversy decomposition}

Let
$pc_{\text{prop}} \in [0,1]$ denote propositional controversy and
$pc_{\text{prag}} \in [0,1]$ pragmatic controversy. Overall perceived
controversy is defined as:
\begin{equation}
pc \,=\, PC(pc_{\text{prop}}, pc_{\text{prag}})
\,=\,
pc_{\text{prop}} \cdot pc_{\text{prag}}.
\end{equation}

This multiplicative form captures the idea that a proposition can be globally
contentious while being relatively uncontroversial in the speaker's local
community, and vice versa.

\subsection{Base utility}

Let $g \in [0,1]$ denote speaker goal strength. The speaker's base context
utility is:
\begin{equation}
U(pc,g)
\,=\,
- w_{pc}\,pc
+ w_g\,g
- w_{\mathrm{int}}\,pc\,g,
\label{eq:model-utility}
\end{equation}
where $w_{pc}, w_g, w_{\mathrm{int}} > 0$.

The first term penalizes controversy, the second rewards persuasive effort,
and the interaction term captures that pursuing a strong goal is increasingly
costly when controversy is high.

\subsection{Noisy threshold intervals}

Each marker $u_i$ is associated with a lower threshold $\theta_i$ and an upper
threshold $\theta_{i+1}$, where the upper threshold is the lower threshold of
the next stronger marker. For the strongest marker, the upper boundary is
$+\infty$.

Thresholds are treated as noisy rather than fixed:
\begin{equation}
\theta_i \sim \mathcal{N}(\mu_i, \sigma_i^2),
\end{equation}
with ordered means
\[
\mu_{\textit{soviel ich wei\ss}}
\;<\;
\mu_{\textit{ja}}
\;<\;
\mu_{\textit{bekanntlich}}.
\]

A marker fits the context to the extent that the utility falls inside its
licensed interval. The corresponding noisy interval-fit term is:
\begin{equation}
\mathrm{Fit}_{\mathrm{interval}}(u_i \mid U)
\;\propto\;
\mathbb{E}_{\theta_i}\left[\sigma\!\left(\lambda(U-\theta_i)\right)\right]
\cdot
\mathbb{E}_{\theta_{i+1}}\left[\sigma\!\left(-\lambda(U-\theta_{i+1})\right)\right],
\label{eq:model-fit}
\end{equation}
where $\sigma(x)=(1+e^{-x})^{-1}$ is the logistic function and $\lambda > 0$
controls boundary sharpness.

The first factor rewards utilities above the lower threshold; the second
factor penalizes utilities above the next stronger marker's lower threshold.
Thus, each marker is associated with an interval rather than a single cutoff.

\subsection{Marker costs and sufficiency blocking}

In addition to interval fit, each marker carries a marker-specific pragmatic
cost $c_u \geq 0$. The weakest marker can be treated as the reference level,
with stronger markers incurring larger costs.

Speaker choice is therefore modeled as:
\begin{equation}
P(u \mid pc,g)
\;\propto\;
\mathrm{Fit}_{\mathrm{interval}}\!\left(u \mid U(pc,g)\right)
\cdot e^{-c_u}.
\label{eq:model-choice}
\end{equation}

This cost term introduces a \emph{weakest-sufficient} bias. A stronger marker
may be licensed by the context, but still lose to a weaker marker if the weaker
option already serves the speaker's purpose sufficiently well. This yields the
prediction of \textbf{sufficiency blocking}: in low-controversy contexts, even
under relatively high goal strength, \textit{bekanntlich} can be blocked if
\textit{ja} or even \textit{soviel ich wei\ss} already provides enough force
at lower cost.

\subsection{Listener inversion}

Given an observed marker $u$, the listener inverts the speaker model to infer
latent context variables:
\begin{equation}
P(pc_{\text{prop}}, pc_{\text{prag}}, g \mid u)
\;\propto\;
P\!\left(u \mid PC(pc_{\text{prop}}, pc_{\text{prag}}), g\right)
\,P(pc_{\text{prop}})\,P(pc_{\text{prag}})\,P(g).
\label{eq:model-posterior}
\end{equation}

Because stronger markers generally require higher utility, they shift posterior
mass toward lower controversy and higher goal strength. However, the cost term
weakens a purely monotonic mapping from marker strength to speaker goal:
observing a weaker marker can reflect either lower utility or the fact that a
stronger marker would have been pragmatically excessive.

\subsection{Adoption model}

The listener's downstream adoption decision depends on the posterior over goal
strength and controversy:
\begin{equation}
\text{logit}\left(P_{\mathrm{adopt}}\right)
\;=
\eta_0
+ \eta_g\,\mathbb{E}[g \mid u,pc]
- \eta_{pc}\,pc
- \eta_{\mathrm{int}}\,\mathbb{E}[g \mid u,pc]\,pc.
\label{eq:model-adopt}
\end{equation}

Here, $\eta_0$ is a baseline adoption tendency, $\eta_g$ captures the positive
effect of inferred speaker goal, $\eta_{pc}$ penalizes controversy, and
$\eta_{\mathrm{int}}$ allows controversy to attenuate the effect of inferred
goal strength.

\subsection{Qualitative predictions}

The model yields the following core predictions:
\begin{enumerate}
  \item \textbf{Goal effect:} holding controversy fixed, higher $g$ increases
        the probability of stronger markers.
  \item \textbf{Controversy effect:} holding $g$ fixed, higher $pc$ increases
        the probability of weaker markers.
  \item \textbf{Attenuated goal effect under controversy:} the impact of $g$
        on marker strength is reduced when $pc$ is high.
  \item \textbf{Sufficiency blocking:} in low-$pc$ regions, a stronger marker
        can be blocked by a weaker marker if the weaker one already does enough.
  \item \textbf{Listener-side inversion:} stronger markers tend to induce lower
        inferred controversy and higher inferred goal strength.
  \item \textbf{Credibility discounting:} even when a strong marker supports
        high inferred goal strength, its effect on action adoption is reduced
        in highly controversial contexts.
\end{enumerate}

\subsection{Interpretive summary}

Relative to a vanilla RSA model that directly softmaxes over utterance utility,
this model treats markers as corresponding to \emph{intervals} on a latent
utility scale. The interval-fit component answers the question:
\enquote{How well does this marker match the current utility region?}
The cost term answers the question:
\enquote{Is this marker stronger or more marked than the speaker needs?}
Only their combination determines actual choice.

\end{document}
